\documentclass{article}
\usepackage{amssymb}
\title{The Native Path: A Working Test Document}
\author{RushTeX Project}
\date{24 October 2026}

\begin{document}
\maketitle

\begin{abstract}
This document is the single growing test case for the native, pure-Rust LaTeX path of
RushTeX. Each section exercises one family of constructs: running text and font
changes, lists, displayed text, tables, floats, mathematics and bibliography. The
document must compile on the native path and produce output identical to pdfLaTeX.
\end{abstract}

\tableofcontents

\section{Running text}\label{sec:text}

Typesetting is the business of turning a stream of characters into a page. A document
author writes paragraphs of plain prose, and the typesetter decides where lines break,
how much space separates words, and where the page ends. The quality of the result
depends on many small decisions, most of which the author never sees. This section
collects the inline devices that an author reaches for while writing prose: emphasis,
font families, font shapes, accents and punctuation.

Emphasis is usually marked with \emph{italic} text, which in an upright paragraph
reverses to roman when nested inside an italic context. Strong statements can be set in
\textbf{bold}, while \textit{explicit italic} and \textsc{small capitals} are useful for
names and terms that deserve a visual distinction without strong emphasis. Source code
and file names belong in \texttt{typewriter type}, and interface labels read best in
\textsf{sans serif}.

Font declarations act on the rest of the current group. The sentence {\bfseries this
phrase is set in bold by a declaration}, and {\itshape this one by an italic
declaration}, so that the surrounding text is unaffected. The older forms \bf and \it
work the same way inside braces, as in {\bf old-style bold} and {\it old-style italic}.

Accents are part of ordinary text. The word caf\'e takes an acute accent, the name
G\"odel takes an umlaut, and the phrase na\"\i ve reads as it should. Quotations are
typed as ``double quotes'' and the apostrophe in don't is a plain character, while
the dashes -- en and --- em -- have their own meanings.

Sentences that trail away end with \ldots{} and the typesetter spaces them as a unit.
The name of the system is written as \TeX{} in the running text and \LaTeX{} in the
document class that builds on it. Non-breaking spaces such as Figure~1 and Section~2
keep a label with its number, and a footnote gives a remark that would interrupt the
argument\footnote{A footnote is set at the foot of the page where its reference
appears, in a smaller size, and it is numbered in sequence with the other footnotes of
the document.}.

Explicit line breaks are useful in addresses and short verse:\\
the first line ends at the double backslash,\\
and the second line starts afresh.

Horizontal and vertical space can be inserted by hand. A paragraph that begins with
\noindent is not indented, and an explicit horizontal gap of \hspace{1em} separates
the two words in this sentence. Vertical space is more common between blocks, so the
next paragraph follows after \vspace{6pt} of extra room.

\noindent The paragraph above was started with the command \texttt{noindent}, so it
has no first-line indentation. The paragraphs that follow are indented as usual, which
is the default behaviour of the article class.

Ordinary paragraphs are separated by a blank line. The line breaking algorithm looks at
the whole paragraph before it commits to any break, which is why a change in one word
can move words several lines further down. The author rarely needs to intervene, but
the mechanisms exist, and the rest of this document uses them where they matter.

\subsection{Font declarations in context}\label{sec:fonts}

A declaration such as \textsf{bfseries} changes the font until the end of the enclosing
group. The next paragraph shows how groups limit the effect of such declarations, and
how the text returns to its original shape when the group closes.

{\bfseries Inside this group the text is bold, and a word such as \emph{this} still
changes to italic within bold text.} Outside the group the text returns to roman, as
the reader can see in this sentence. The same holds for {\itshape the slanted shape}
and for {\sffamily the sans serif family}, which all end where their braces end.

\subsection{Spacing and hyphenation}\label{sec:spacing}

Word spacing in a justified line is not constant. The stretch that a line receives is
shared among its interword glue, and the amount of stretch depends on how much room the
line has after its words are placed. A line with few words and a large gap is
noticeably loose, and a line with many words and little room may be tight. The
typesetter tries to keep all lines of a paragraph as even as possible, and it accepts
a slightly loose line when the alternative is a worse paragraph as a whole.

Long words that do not fit at the end of a line are hyphenated at points chosen from a
table of patterns. The table is built from word lists and the rules that govern where
a word may be divided, and it lets the typesetter break words such as typesetting,
hyphenation, algorithm and reproducibility at reasonable places. A hyphen at the end of
a line is a cost, and two such hyphens in a row are a larger cost, so the algorithm
avoids them where it can.

Some words must never be divided, and the author can say so by writing the word inside
a group that forbids breaks, or by joining the parts of a name with a tie. The name
Ludwig~van Beethoven keeps its two parts together, and a unit such as 12~pt does the
same for the number and the unit.

Sentences that end with a capital letter, such as the initials in J.~R.\ Smith, have a
different space after them, and the author must mark an end of sentence when the
capital is not the end of a sentence. The rule is simple, but it is easy to forget, and
the document below uses a backslash space in one case to show how it is written.

\section{Lists}\label{sec:lists}

Lists organise material that has no natural sentence structure. The three list
environments of the article class differ in how they mark their items and how they
number them. The first is the unordered list, which uses a bullet for each item.

\begin{itemize}
  \item The first item of a bulleted list is a plain sentence.
  \item The second item is longer and contains an \emph{emphasised} word, so that the
        wrapped line is indented to the same position as the text of the item.
  \item A third item may itself hold a sublist:
    \begin{itemize}
      \item the nested list uses a dash as its marker,
      \item and its items are indented further than the parent items.
    \end{itemize}
\end{itemize}

Numbered lists are used when the order of the items matters, for example in a
procedure. The following enumeration is nested two levels deep, and its first item
carries a label so that the text can refer to it later. The reference to that item
appears in Section~\ref{sec:tables}, which is written after this list.

\begin{enumerate}
  \item\label{item:first} Obtain the source files of the document and the packages that
        it loads.
  \item Compile the document once on the native path, which reads the auxiliary
        information of the previous run, if there is one.
        \begin{enumerate}
          \item The first sub-item describes a check that the author can make by hand.
          \item The second sub-item describes what happens when the check fails.
        \end{enumerate}
  \item Compare the result against the reference output and report any page that
        differs.
\end{enumerate}

The first step of the procedure is item~\ref{item:first} in the list above, and the
reader who wants the reasons behind the procedure will find them in
Section~\ref{sec:displays}.

Definition lists give a term followed by its explanation. The description environment
sets each term in bold and indents the explanation beneath it.

\begin{description}
  \item[Kerning] An adjustment of the space between two particular letters, such as
        the pair A and V, so that the shapes look evenly spaced.
  \item[Ligature] A single glyph that replaces a sequence of letters, such as the
        fi ligature in the word first.
  \item[Leading] The distance from one baseline to the next, measured in points.
\end{description}

The three kinds of list above are the core of most structured writing. An author who
needs a list inside a list simply nests the environment, and each level takes its own
label and its own indentation. The nesting depth in the article class is limited only
by the page, although a list nested more than three levels deep is usually a sign that
the material wants a table or a section instead.

Lists also interact with the page. A list is a sequence of paragraphs with extra space
between items, and the page builder may break the page between two items, as it may
between two paragraphs. The space at the top and bottom of a list is controlled by the
document class, so that a list opens and closes with the same visual distance as a
displayed equation.

\section{Displayed text}\label{sec:displays}

Some text is set apart from the body of the document to show that it is a quotation, a
centred title, or a block of literal material. Each environment has its own margins and
its own vertical spacing, and the environments below show the common cases.

\begin{quote}
A quotation environment sets its text with extra margin on both sides and without
paragraph indentation. It is intended for a short passage that the author wishes to
present as a separate block of prose, without the extra emphasis of a block quotation.
\end{quote}

\begin{quotation}
The quotation environment differs from the quote environment in that it indents the
first line of each of its paragraphs. A long passage that runs over several paragraphs
therefore keeps its internal structure, as a reader expects from a transcript of a
speech or a letter. This quotation has two paragraphs, so the second one can be seen to
start with an indent.

The second paragraph of the quotation shows the indentation of a continuing paragraph.
\end{quotation}

\begin{center}
Centred text on a line of its own\\
keeps each of its lines in the middle of the text block.
\end{center}

\begin{flushleft}
Flush left text keeps the left margin straight and lets the right margin
ragged, which suits material such as addresses and minutes of a meeting.
\end{flushleft}

\begin{flushright}
Flush right text keeps the right margin straight, which is seldom useful for
prose but is useful for a signature line.
\end{flushright}

Verbatim text is typeset exactly as it is typed, with spaces, line breaks and special
characters such as the percent sign and the braces kept as they stand:

\begin{verbatim}
  for i in 1..n {
      total += a[i] * b[i];   % dot product, not a comment
  }
\end{verbatim}

The following block uses a centring declaration inside a group, which keeps the
alignment local to the braces, and then a ragged right declaration on the next block.

{\centering A centred block made with a declaration, which ends at the closing
brace and leaves the surrounding text alone.\par}

{\raggedright A block whose lines are not stretched to the full measure, so the
right margin is ragged. The declaration ends with the group, and the next paragraph
returns to the usual justified setting of the text.\par}

Paragraphs of ordinary text after these blocks are justified again, as the reader can
see here.

\section{Tables}\label{sec:tables}

Tables set data in rows and columns. The tabular environment takes a column
specification, and each column is either left, centred or right aligned. Vertical bars
draw rules between columns. Section~\ref{sec:displays} explains the centring that surrounds each table in
this section.

\subsection{Basic alignment and rules}\label{sec:tables-basic}

\begin{center}
\begin{tabular}{l c r}
Method & Passes & Time (ms) \\
Greedy fill      & 1 & 0.4 \\
Knuth--Plass     & 2 & 3.1 \\
Optimal fit      & 3 & 4.7 \\
\end{tabular}
\end{center}

\begin{center}
\begin{tabular}{|l|c|r|}
Language & Typing & Speed \\
C        & static  & fast   \\
Python   & dynamic & slow   \\
Rust     & static  & fast   \\
\end{tabular}
\end{center}

\subsection{Padding and repeated columns}\label{sec:tables-advanced}

The column specification may fix the padding around a column with an empty at-sign
argument, which removes the space at one end of the table. Repeated columns are written
with a star argument, which expands to a run of identical columns.

\begin{center}
\begin{tabular}{@{}*{3}{c}@{}}
Left & Middle & Right \\
1 & 2 & 3 \\
4 & 5 & 6 \\
\end{tabular}
\end{center}

Each column keeps its own alignment in every row, so the numbers in the table below
line up on their right edges whatever the width of their labels.

\begin{center}
\begin{tabular}{l c r}
Run & Pages & Seconds \\
Alpha & 12 & 0.8 \\
Beta  & 29 & 1.9 \\
Gamma & 61 & 4.2 \\
\end{tabular}
\end{center}

Extra vertical space between rows is written after a row break with an optional length,
as in \texttt{[2pt]}. The table below uses that spacing to separate groups of rows.

\begin{center}
\begin{tabular}{l r}
Operation & Cost \\[2pt]
Read      & 3 \\
Write     & 7 \\[2pt]
Reference & 11 \\
\end{tabular}
\end{center}

Tables with many rows are common in reports, and the author who sets them by hand often
spends more time on alignment than on content. The tabular environment reduces that
work to a column specification and a list of rows. Each row ends with a double
backslash, and each cell is separated from its neighbour by an ampersand. The
alignment of each column is taken from the specification, so a change of one letter
moves the whole column without any other edit.

The tables so far show the alignment letters, the padding that the at-sign notation
controls and the repetition of a column pattern. Horizontal rules, partial rules and
cells that span several columns are set by the same machinery.

\subsection{Rules and spanning cells}\label{sec:tables-rules}
Table rules come in three kinds: a full rule across every column, a partial rule
across some of them, and a doubled rule made of two full rules in a row. A heading can
span all the columns of the table, and a single cell can span two of them.

\begin{center}
\begin{tabular}{|l|c|r|}
\hline
\multicolumn{3}{|c|}{\bfseries Measured costs} \\
\hline\hline
Phase & Calls & Time \\
\cline{2-3}
Scan    & 120 & 3.5 \\
Expand  & 940 & 12.0 \\
\multicolumn{2}{|l|}{Total, both phases} & 15.5 \\
\hline
\end{tabular}
\end{center}

\section{Floats}\label{sec:floats}

Floating material such as tables and figures is moved by the page builder to a position
where it fits. The placement specifier suggests where the float may go: \texttt{h} for
here, \texttt{t} for the top of a page, \texttt{b} for the bottom and \texttt{p} for a
page of floats. The author can give several letters at once, as in \texttt{[htbp]}, and
the page builder tries them in turn.

Table~\ref{tab:results} appears before the text that discusses it, so the reader meets
the reference first and the float later. The float is placed here by its own
specifier.

\begin{table}[h]
\centering
\begin{tabular}{l c r}
Input size & Lines & Pages \\
Short      & 40    & 1     \\
Medium     & 400   & 12    \\
Long       & 4000  & 120   \\
\end{tabular}
\caption{Pages produced for three input sizes.}\label{tab:results}
\end{table}

The results in Table~\ref{tab:results} grow with the input, and the page counts follow
the line counts closely. The next float is a figure with a caption long enough to wrap
onto a second line, and it is placed at the top of the page by its specifier.

\begin{figure}[t]
\centering
\mbox{\textsf{Input}\quad$\rightarrow$\quad\texttt{Parser}\quad$\rightarrow$\quad\textsf{Layout}\quad$\rightarrow$\quad\textsf{PDF}}
\caption{The stages of the native pipeline, from the characters of the source file to the
bytes of the output file, shown as a chain of boxes in the order in which the data flows
through them.}\label{fig:pipeline}
\end{figure}

Figure~\ref{fig:pipeline} shows the chain of stages that the native path runs in a
single pass. The figure caption is set as a paragraph because it runs over two lines.
A table placed at the bottom of a page is specified with \texttt{[b]}, as the next
example shows.

\begin{table}[b]
\centering
\begin{tabular}{c c}
Feature & State \\
Lists        & done \\
Tables       & done \\
\end{tabular}
\caption{Feature status at the time of writing.}
\end{table}

A float with both the top and the bottom permitted is written \texttt{[tb]}, and a float
that may go anywhere in the sequence is written \texttt{[htbp]}. The next figure uses
the latter specifier and a label that the text refers to from a later section.

\begin{figure}[htbp]
\centering
\mbox{\texttt{figure}} \quad \mbox{Text placed in a box}
\caption{A short caption.}\label{fig:box}
\end{figure}

Figure~\ref{fig:box} is referred to here, after the float itself. The references to the
tables and figures of this section are resolved in the second pass, which is the reason
that the reader sees the correct numbers even for references that come before their
floats.

A float is a block of material that the page builder may move away from the place where
it appears in the source. Its purpose is to keep large items, such as a table of
results or a diagram, from splitting across a page break or from leaving a large gap
at the foot of a page. The page builder reads the placement specifier of the float and
considers each candidate position in turn. If no candidate fits, the float waits for
the next page, and the material that follows it continues to fill the current page.

A caption belongs to its float and is numbered with it. The number of a table and the
number of a figure are kept in separate counters, so that a document may contain a
table 1 and a figure 1 without conflict. The label placed after the caption records the
number, and the reference in the text picks up the number from the label when the
document is compiled. This is the reason that a reference may come before its float in
the source, as the references to the tables above do.

\section{Mathematics}\label{sec:math}

Mathematical notation is set in math mode, which uses its own fonts and spacing rules.
Inline formulae are written between dollar signs. The quantity $\frac{a+b}{c-d}$ is a
fraction, and $\sqrt{x^2+1}$ is a square root. A superscript such as $x^{2}$ and a
subscript such as $a_{i}$ are common, and both can appear together, as in $x_{i}^{2}$.

Greek letters are typed by name, so $\alpha$, $\beta$ and $\gamma$ are available, and so
is the capital $\Gamma$. The sum $\sum_{k=1}^{n} k$ and the integral $\int_0^1 x\,dx$
appear inline with their limits set beside them, as is customary in running text. The
limit $\lim_{n\to\infty} \frac{1}{n} = 0$ also appears inline.

Sets are typeset with blackboard and calligraphic letters. The set of real numbers is
$\mathbb{R}$, and a function that bounds another is written $f(n) \in \mathcal{O}(n^2)$.
A vector is written in bold, as $\mathbf{v}$, and its length is $\|\mathbf{v}\|$. The
parentheses in $\left( \frac{x}{y} \right)$ grow with the fraction they enclose.

Displayed formulae are set on their own line, centred, with a larger style than inline
formulae. The formula below gives the sum of the first $n$ integers.
\[
  \sum_{k=1}^{n} k = \frac{n(n+1)}{2}
\]

Numbered equations carry a label, so that the text can refer to them by number. The
first equation of this section is the definition of the Gaussian integral, and the
second is the quadratic formula, which is used in Section~\ref{sec:lists}.

\begin{equation}\label{eq:gauss}
  \int_{-\infty}^{\infty} e^{-x^{2}}\,dx = \sqrt{\pi}
\end{equation}

\begin{equation}\label{eq:quadratic}
  x = \frac{-b \pm \sqrt{b^{2}-4ac}}{2a}
\end{equation}

Equation~\ref{eq:gauss} gives the Gaussian integral, and
Equation~\ref{eq:quadratic} gives the roots of a quadratic polynomial. The coefficients
$a$, $b$ and $c$ in Equation~\ref{eq:quadratic} are real numbers with $a \neq 0$.

A last display shows a sum over a product, with a subscript that is itself a formula:
\[
  \prod_{i=1}^{m} \left( 1 + \frac{1}{i^{2}} \right) \leq e^{\pi^{2}/6}
\]

Mathematical formulae are read as a sequence of symbols with a structure, and the layout
rules give each symbol a size and a position. Inline formulae are set in the text style,
which keeps the lines of a paragraph at a uniform height, while displayed formulae are
set in the display style, which allows larger operators and fractions with larger
numerators and denominators. The reader sees the difference in the sums and products
above, where the limits sit above and below the symbol in the display, and beside it in
the running text.

The quadratic formula in Equation~\ref{eq:quadratic} is a familiar example of a
fraction with a root in its numerator. The square root of the discriminant
$b^{2}-4ac$ is set with a radical sign that grows with its contents, and the
fraction line extends beneath the whole numerator. The parentheses of the next example
are sized in the same way, so that the reader can see the structure at a glance.

\subsection{Displayed derivations}\label{sec:derivations}

A derivation shows a chain of equalities, each following from the one before. Each line
of the chain is written as a separate display, and the text between the displays gives
the reason for each step. The first derivation below expands the square of a binomial,
and the second uses the result to evaluate a finite sum of squares.
\[
  (a+b)^{2} = a^{2} + 2ab + b^{2}
\]

Applying the same expansion to the terms of a sum gives a closed form for the sum of the
squares of the first $n$ integers, which is the second formula of this subsection:
\[
  \sum_{k=1}^{n} k^{2} = \frac{n(n+1)(2n+1)}{6}
\]

The formula holds for every positive integer $n$, and it can be proved by induction on
$n$, as the reader may check for the first few values. Each step of the induction
uses the identity above and the fact that the sum for $n+1$ is the sum for $n$ with one
more term added.

The third display is a limit that defines the constant $e$. It is set on its own line
because it is a result that the text refers to, and the reader may wish to compare it
with the bound given in the last display of the previous section:
\[
  e = \lim_{n\to\infty} \left( 1 + \frac{1}{n} \right)^{n}
\]

Derivations that are long are best broken into several displays, each short enough to
fit on one line of the text block, so that the reader can follow the argument without
turning the page.

\section{Bibliography and citations}\label{sec:refs}

The native path follows the algorithms of the typesetting system closely, and the
descriptions of those algorithms are found in the standard references. The line
breaking method used by the paragraph builder is described by Knuth and
Plass~\cite{knuth81}, while the page builder and the float placement are described in
the reference manual~\cite{knuth86}. The layout of mathematical formulae follows the
rules of the appendix on mathematics in that same manual. The chain of stages in
Figure~\ref{fig:pipeline} is a common design for compilers in general, as Aho and
colleagues describe~\cite{aho06}.

\begin{thebibliography}{9}
\bibitem{knuth81}
  Donald E.~Knuth and Michael F.~Plass.
  \emph{Breaking paragraphs into lines}.
  Software: Practice and Experience, 11(11), 1981.

\bibitem{knuth86}
  Donald E.~Knuth.
  \emph{The \TeX book}.
  Addison-Wesley, 1986.

\bibitem{aho06}
  Alfred V.~Aho, Monica S.~Lam, Ravi Sethi and Jeffrey D.~Ullman.
  \emph{Compilers: Principles, Techniques, and Tools}.
  Addison-Wesley, second edition, 2006.
\end{thebibliography}



\section{Running text}\label{sec:text-1}

Typesetting is the business of turning a stream of characters into a page. A document
author writes paragraphs of plain prose, and the typesetter decides where lines break,
how much space separates words, and where the page ends. The quality of the result
depends on many small decisions, most of which the author never sees. This section
collects the inline devices that an author reaches for while writing prose: emphasis,
font families, font shapes, accents and punctuation.

Emphasis is usually marked with \emph{italic} text, which in an upright paragraph
reverses to roman when nested inside an italic context. Strong statements can be set in
\textbf{bold}, while \textit{explicit italic} and \textsc{small capitals} are useful for
names and terms that deserve a visual distinction without strong emphasis. Source code
and file names belong in \texttt{typewriter type}, and interface labels read best in
\textsf{sans serif}.

Font declarations act on the rest of the current group. The sentence {\bfseries this
phrase is set in bold by a declaration}, and {\itshape this one by an italic
declaration}, so that the surrounding text is unaffected. The older forms \bf and \it
work the same way inside braces, as in {\bf old-style bold} and {\it old-style italic}.

Accents are part of ordinary text. The word caf\'e takes an acute accent, the name
G\"odel takes an umlaut, and the phrase na\"\i ve reads as it should. Quotations are
typed as ``double quotes'' and the apostrophe in don't is a plain character, while
the dashes -- en and --- em -- have their own meanings.

Sentences that trail away end with \ldots{} and the typesetter spaces them as a unit.
The name of the system is written as \TeX{} in the running text and \LaTeX{} in the
document class that builds on it. Non-breaking spaces such as Figure~1 and Section~2
keep a label with its number, and a footnote gives a remark that would interrupt the
argument\footnote{A footnote is set at the foot of the page where its reference
appears, in a smaller size, and it is numbered in sequence with the other footnotes of
the document.}.

Explicit line breaks are useful in addresses and short verse:\\
the first line ends at the double backslash,\\
and the second line starts afresh.

Horizontal and vertical space can be inserted by hand. A paragraph that begins with
\noindent is not indented, and an explicit horizontal gap of \hspace{1em} separates
the two words in this sentence. Vertical space is more common between blocks, so the
next paragraph follows after \vspace{6pt} of extra room.

\noindent The paragraph above was started with the command \texttt{noindent}, so it
has no first-line indentation. The paragraphs that follow are indented as usual, which
is the default behaviour of the article class.

Ordinary paragraphs are separated by a blank line. The line breaking algorithm looks at
the whole paragraph before it commits to any break, which is why a change in one word
can move words several lines further down. The author rarely needs to intervene, but
the mechanisms exist, and the rest of this document uses them where they matter.

\subsection{Font declarations in context}\label{sec:fonts-1}

A declaration such as \textsf{bfseries} changes the font until the end of the enclosing
group. The next paragraph shows how groups limit the effect of such declarations, and
how the text returns to its original shape when the group closes.

{\bfseries Inside this group the text is bold, and a word such as \emph{this} still
changes to italic within bold text.} Outside the group the text returns to roman, as
the reader can see in this sentence. The same holds for {\itshape the slanted shape}
and for {\sffamily the sans serif family}, which all end where their braces end.

\subsection{Spacing and hyphenation}\label{sec:spacing-1}

Word spacing in a justified line is not constant. The stretch that a line receives is
shared among its interword glue, and the amount of stretch depends on how much room the
line has after its words are placed. A line with few words and a large gap is
noticeably loose, and a line with many words and little room may be tight. The
typesetter tries to keep all lines of a paragraph as even as possible, and it accepts
a slightly loose line when the alternative is a worse paragraph as a whole.

Long words that do not fit at the end of a line are hyphenated at points chosen from a
table of patterns. The table is built from word lists and the rules that govern where
a word may be divided, and it lets the typesetter break words such as typesetting,
hyphenation, algorithm and reproducibility at reasonable places. A hyphen at the end of
a line is a cost, and two such hyphens in a row are a larger cost, so the algorithm
avoids them where it can.

Some words must never be divided, and the author can say so by writing the word inside
a group that forbids breaks, or by joining the parts of a name with a tie. The name
Ludwig~van Beethoven keeps its two parts together, and a unit such as 12~pt does the
same for the number and the unit.

Sentences that end with a capital letter, such as the initials in J.~R.\ Smith, have a
different space after them, and the author must mark an end of sentence when the
capital is not the end of a sentence. The rule is simple, but it is easy to forget, and
the document below uses a backslash space in one case to show how it is written.

\section{Lists}\label{sec:lists-1}

Lists organise material that has no natural sentence structure. The three list
environments of the article class differ in how they mark their items and how they
number them. The first is the unordered list, which uses a bullet for each item.

\begin{itemize}
  \item The first item of a bulleted list is a plain sentence.
  \item The second item is longer and contains an \emph{emphasised} word, so that the
        wrapped line is indented to the same position as the text of the item.
  \item A third item may itself hold a sublist:
    \begin{itemize}
      \item the nested list uses a dash as its marker,
      \item and its items are indented further than the parent items.
    \end{itemize}
\end{itemize}

Numbered lists are used when the order of the items matters, for example in a
procedure. The following enumeration is nested two levels deep, and its first item
carries a label so that the text can refer to it later. The reference to that item
appears in Section~\ref{sec:tables-1}, which is written after this list.

\begin{enumerate}
  \item\label{item:first-1} Obtain the source files of the document and the packages that
        it loads.
  \item Compile the document once on the native path, which reads the auxiliary
        information of the previous run, if there is one.
        \begin{enumerate}
          \item The first sub-item describes a check that the author can make by hand.
          \item The second sub-item describes what happens when the check fails.
        \end{enumerate}
  \item Compare the result against the reference output and report any page that
        differs.
\end{enumerate}

The first step of the procedure is item~\ref{item:first-1} in the list above, and the
reader who wants the reasons behind the procedure will find them in
Section~\ref{sec:displays-1}.

Definition lists give a term followed by its explanation. The description environment
sets each term in bold and indents the explanation beneath it.

\begin{description}
  \item[Kerning] An adjustment of the space between two particular letters, such as
        the pair A and V, so that the shapes look evenly spaced.
  \item[Ligature] A single glyph that replaces a sequence of letters, such as the
        fi ligature in the word first.
  \item[Leading] The distance from one baseline to the next, measured in points.
\end{description}

The three kinds of list above are the core of most structured writing. An author who
needs a list inside a list simply nests the environment, and each level takes its own
label and its own indentation. The nesting depth in the article class is limited only
by the page, although a list nested more than three levels deep is usually a sign that
the material wants a table or a section instead.

Lists also interact with the page. A list is a sequence of paragraphs with extra space
between items, and the page builder may break the page between two items, as it may
between two paragraphs. The space at the top and bottom of a list is controlled by the
document class, so that a list opens and closes with the same visual distance as a
displayed equation.

\section{Displayed text}\label{sec:displays-1}

Some text is set apart from the body of the document to show that it is a quotation, a
centred title, or a block of literal material. Each environment has its own margins and
its own vertical spacing, and the environments below show the common cases.

\begin{quote}
A quotation environment sets its text with extra margin on both sides and without
paragraph indentation. It is intended for a short passage that the author wishes to
present as a separate block of prose, without the extra emphasis of a block quotation.
\end{quote}

\begin{quotation}
The quotation environment differs from the quote environment in that it indents the
first line of each of its paragraphs. A long passage that runs over several paragraphs
therefore keeps its internal structure, as a reader expects from a transcript of a
speech or a letter. This quotation has two paragraphs, so the second one can be seen to
start with an indent.

The second paragraph of the quotation shows the indentation of a continuing paragraph.
\end{quotation}

\begin{center}
Centred text on a line of its own\\
keeps each of its lines in the middle of the text block.
\end{center}

\begin{flushleft}
Flush left text keeps the left margin straight and lets the right margin
ragged, which suits material such as addresses and minutes of a meeting.
\end{flushleft}

\begin{flushright}
Flush right text keeps the right margin straight, which is seldom useful for
prose but is useful for a signature line.
\end{flushright}

Verbatim text is typeset exactly as it is typed, with spaces, line breaks and special
characters such as the percent sign and the braces kept as they stand:

\begin{verbatim}
  for i in 1..n {
      total += a[i] * b[i];   % dot product, not a comment
  }
\end{verbatim}

The following block uses a centring declaration inside a group, which keeps the
alignment local to the braces, and then a ragged right declaration on the next block.

{\centering A centred block made with a declaration, which ends at the closing
brace and leaves the surrounding text alone.\par}

{\raggedright A block whose lines are not stretched to the full measure, so the
right margin is ragged. The declaration ends with the group, and the next paragraph
returns to the usual justified setting of the text.\par}

Paragraphs of ordinary text after these blocks are justified again, as the reader can
see here.

\section{Tables}\label{sec:tables-1}

Tables set data in rows and columns. The tabular environment takes a column
specification, and each column is either left, centred or right aligned. Vertical bars
draw rules between columns. Section~\ref{sec:displays-1} explains the centring that surrounds each table in
this section.

\subsection{Basic alignment and rules}\label{sec:tables-basic-1}

\begin{center}
\begin{tabular}{l c r}
Method & Passes & Time (ms) \\
Greedy fill      & 1 & 0.4 \\
Knuth--Plass     & 2 & 3.1 \\
Optimal fit      & 3 & 4.7 \\
\end{tabular}
\end{center}

\begin{center}
\begin{tabular}{|l|c|r|}
Language & Typing & Speed \\
C        & static  & fast   \\
Python   & dynamic & slow   \\
Rust     & static  & fast   \\
\end{tabular}
\end{center}

\subsection{Padding and repeated columns}\label{sec:tables-advanced-1}

The column specification may fix the padding around a column with an empty at-sign
argument, which removes the space at one end of the table. Repeated columns are written
with a star argument, which expands to a run of identical columns.

\begin{center}
\begin{tabular}{@{}*{3}{c}@{}}
Left & Middle & Right \\
1 & 2 & 3 \\
4 & 5 & 6 \\
\end{tabular}
\end{center}

Each column keeps its own alignment in every row, so the numbers in the table below
line up on their right edges whatever the width of their labels.

\begin{center}
\begin{tabular}{l c r}
Run & Pages & Seconds \\
Alpha & 12 & 0.8 \\
Beta  & 29 & 1.9 \\
Gamma & 61 & 4.2 \\
\end{tabular}
\end{center}

Extra vertical space between rows is written after a row break with an optional length,
as in \texttt{[2pt]}. The table below uses that spacing to separate groups of rows.

\begin{center}
\begin{tabular}{l r}
Operation & Cost \\[2pt]
Read      & 3 \\
Write     & 7 \\[2pt]
Reference & 11 \\
\end{tabular}
\end{center}

Tables with many rows are common in reports, and the author who sets them by hand often
spends more time on alignment than on content. The tabular environment reduces that
work to a column specification and a list of rows. Each row ends with a double
backslash, and each cell is separated from its neighbour by an ampersand. The
alignment of each column is taken from the specification, so a change of one letter
moves the whole column without any other edit.

The tables so far show the alignment letters, the padding that the at-sign notation
controls and the repetition of a column pattern. Horizontal rules, partial rules and
cells that span several columns are set by the same machinery.

\subsection{Rules and spanning cells}\label{sec:tables-rules-1}
Table rules come in three kinds: a full rule across every column, a partial rule
across some of them, and a doubled rule made of two full rules in a row. A heading can
span all the columns of the table, and a single cell can span two of them.

\begin{center}
\begin{tabular}{|l|c|r|}
\hline
\multicolumn{3}{|c|}{\bfseries Measured costs} \\
\hline\hline
Phase & Calls & Time \\
\cline{2-3}
Scan    & 120 & 3.5 \\
Expand  & 940 & 12.0 \\
\multicolumn{2}{|l|}{Total, both phases} & 15.5 \\
\hline
\end{tabular}
\end{center}

\section{Floats}\label{sec:floats-1}

Floating material such as tables and figures is moved by the page builder to a position
where it fits. The placement specifier suggests where the float may go: \texttt{h} for
here, \texttt{t} for the top of a page, \texttt{b} for the bottom and \texttt{p} for a
page of floats. The author can give several letters at once, as in \texttt{[htbp]}, and
the page builder tries them in turn.

Table~\ref{tab:results-1} appears before the text that discusses it, so the reader meets
the reference first and the float later. The float is placed here by its own
specifier.

\begin{table}[h]
\centering
\begin{tabular}{l c r}
Input size & Lines & Pages \\
Short      & 40    & 1     \\
Medium     & 400   & 12    \\
Long       & 4000  & 120   \\
\end{tabular}
\caption{Pages produced for three input sizes.}\label{tab:results-1}
\end{table}

The results in Table~\ref{tab:results-1} grow with the input, and the page counts follow
the line counts closely. The next float is a figure with a caption long enough to wrap
onto a second line, and it is placed at the top of the page by its specifier.

\begin{figure}[t]
\centering
\mbox{\textsf{Input}\quad$\rightarrow$\quad\texttt{Parser}\quad$\rightarrow$\quad\textsf{Layout}\quad$\rightarrow$\quad\textsf{PDF}}
\caption{The stages of the native pipeline, from the characters of the source file to the
bytes of the output file, shown as a chain of boxes in the order in which the data flows
through them.}\label{fig:pipeline-1}
\end{figure}

Figure~\ref{fig:pipeline-1} shows the chain of stages that the native path runs in a
single pass. The figure caption is set as a paragraph because it runs over two lines.
A table placed at the bottom of a page is specified with \texttt{[b]}, as the next
example shows.

\begin{table}[b]
\centering
\begin{tabular}{c c}
Feature & State \\
Lists        & done \\
Tables       & done \\
\end{tabular}
\caption{Feature status at the time of writing.}
\end{table}

A float with both the top and the bottom permitted is written \texttt{[tb]}, and a float
that may go anywhere in the sequence is written \texttt{[htbp]}. The next figure uses
the latter specifier and a label that the text refers to from a later section.

\begin{figure}[htbp]
\centering
\mbox{\texttt{figure}} \quad \mbox{Text placed in a box}
\caption{A short caption.}\label{fig:box-1}
\end{figure}

Figure~\ref{fig:box-1} is referred to here, after the float itself. The references to the
tables and figures of this section are resolved in the second pass, which is the reason
that the reader sees the correct numbers even for references that come before their
floats.

A float is a block of material that the page builder may move away from the place where
it appears in the source. Its purpose is to keep large items, such as a table of
results or a diagram, from splitting across a page break or from leaving a large gap
at the foot of a page. The page builder reads the placement specifier of the float and
considers each candidate position in turn. If no candidate fits, the float waits for
the next page, and the material that follows it continues to fill the current page.

A caption belongs to its float and is numbered with it. The number of a table and the
number of a figure are kept in separate counters, so that a document may contain a
table 1 and a figure 1 without conflict. The label placed after the caption records the
number, and the reference in the text picks up the number from the label when the
document is compiled. This is the reason that a reference may come before its float in
the source, as the references to the tables above do.

\section{Mathematics}\label{sec:math-1}

Mathematical notation is set in math mode, which uses its own fonts and spacing rules.
Inline formulae are written between dollar signs. The quantity $\frac{a+b}{c-d}$ is a
fraction, and $\sqrt{x^2+1}$ is a square root. A superscript such as $x^{2}$ and a
subscript such as $a_{i}$ are common, and both can appear together, as in $x_{i}^{2}$.

Greek letters are typed by name, so $\alpha$, $\beta$ and $\gamma$ are available, and so
is the capital $\Gamma$. The sum $\sum_{k=1}^{n} k$ and the integral $\int_0^1 x\,dx$
appear inline with their limits set beside them, as is customary in running text. The
limit $\lim_{n\to\infty} \frac{1}{n} = 0$ also appears inline.

Sets are typeset with blackboard and calligraphic letters. The set of real numbers is
$\mathbb{R}$, and a function that bounds another is written $f(n) \in \mathcal{O}(n^2)$.
A vector is written in bold, as $\mathbf{v}$, and its length is $\|\mathbf{v}\|$. The
parentheses in $\left( \frac{x}{y} \right)$ grow with the fraction they enclose.

Displayed formulae are set on their own line, centred, with a larger style than inline
formulae. The formula below gives the sum of the first $n$ integers.
\[
  \sum_{k=1}^{n} k = \frac{n(n+1)}{2}
\]

Numbered equations carry a label, so that the text can refer to them by number. The
first equation of this section is the definition of the Gaussian integral, and the
second is the quadratic formula, which is used in Section~\ref{sec:lists-1}.

\begin{equation}\label{eq:gauss-1}
  \int_{-\infty}^{\infty} e^{-x^{2}}\,dx = \sqrt{\pi}
\end{equation}

\begin{equation}\label{eq:quadratic-1}
  x = \frac{-b \pm \sqrt{b^{2}-4ac}}{2a}
\end{equation}

Equation~\ref{eq:gauss-1} gives the Gaussian integral, and
Equation~\ref{eq:quadratic-1} gives the roots of a quadratic polynomial. The coefficients
$a$, $b$ and $c$ in Equation~\ref{eq:quadratic-1} are real numbers with $a \neq 0$.

A last display shows a sum over a product, with a subscript that is itself a formula:
\[
  \prod_{i=1}^{m} \left( 1 + \frac{1}{i^{2}} \right) \leq e^{\pi^{2}/6}
\]

Mathematical formulae are read as a sequence of symbols with a structure, and the layout
rules give each symbol a size and a position. Inline formulae are set in the text style,
which keeps the lines of a paragraph at a uniform height, while displayed formulae are
set in the display style, which allows larger operators and fractions with larger
numerators and denominators. The reader sees the difference in the sums and products
above, where the limits sit above and below the symbol in the display, and beside it in
the running text.

The quadratic formula in Equation~\ref{eq:quadratic-1} is a familiar example of a
fraction with a root in its numerator. The square root of the discriminant
$b^{2}-4ac$ is set with a radical sign that grows with its contents, and the
fraction line extends beneath the whole numerator. The parentheses of the next example
are sized in the same way, so that the reader can see the structure at a glance.

\subsection{Displayed derivations}\label{sec:derivations-1}

A derivation shows a chain of equalities, each following from the one before. Each line
of the chain is written as a separate display, and the text between the displays gives
the reason for each step. The first derivation below expands the square of a binomial,
and the second uses the result to evaluate a finite sum of squares.
\[
  (a+b)^{2} = a^{2} + 2ab + b^{2}
\]

Applying the same expansion to the terms of a sum gives a closed form for the sum of the
squares of the first $n$ integers, which is the second formula of this subsection:
\[
  \sum_{k=1}^{n} k^{2} = \frac{n(n+1)(2n+1)}{6}
\]

The formula holds for every positive integer $n$, and it can be proved by induction on
$n$, as the reader may check for the first few values. Each step of the induction
uses the identity above and the fact that the sum for $n+1$ is the sum for $n$ with one
more term added.

The third display is a limit that defines the constant $e$. It is set on its own line
because it is a result that the text refers to, and the reader may wish to compare it
with the bound given in the last display of the previous section:
\[
  e = \lim_{n\to\infty} \left( 1 + \frac{1}{n} \right)^{n}
\]

Derivations that are long are best broken into several displays, each short enough to
fit on one line of the text block, so that the reader can follow the argument without
turning the page.

\section{Bibliography and citations}\label{sec:refs-1}

The native path follows the algorithms of the typesetting system closely, and the
descriptions of those algorithms are found in the standard references. The line
breaking method used by the paragraph builder is described by Knuth and
Plass~\cite{knuth81-1}, while the page builder and the float placement are described in
the reference manual~\cite{knuth86-1}. The layout of mathematical formulae follows the
rules of the appendix on mathematics in that same manual. The chain of stages in
Figure~\ref{fig:pipeline-1} is a common design for compilers in general, as Aho and
colleagues describe~\cite{aho06-1}.

\begin{thebibliography}{9}
\bibitem{knuth81-1}
  Donald E.~Knuth and Michael F.~Plass.
  \emph{Breaking paragraphs into lines}.
  Software: Practice and Experience, 11(11), 1981.

\bibitem{knuth86-1}
  Donald E.~Knuth.
  \emph{The \TeX book}.
  Addison-Wesley, 1986.

\bibitem{aho06-1}
  Alfred V.~Aho, Monica S.~Lam, Ravi Sethi and Jeffrey D.~Ullman.
  \emph{Compilers: Principles, Techniques, and Tools}.
  Addison-Wesley, second edition, 2006.
\end{thebibliography}



\section{Running text}\label{sec:text-2}

Typesetting is the business of turning a stream of characters into a page. A document
author writes paragraphs of plain prose, and the typesetter decides where lines break,
how much space separates words, and where the page ends. The quality of the result
depends on many small decisions, most of which the author never sees. This section
collects the inline devices that an author reaches for while writing prose: emphasis,
font families, font shapes, accents and punctuation.

Emphasis is usually marked with \emph{italic} text, which in an upright paragraph
reverses to roman when nested inside an italic context. Strong statements can be set in
\textbf{bold}, while \textit{explicit italic} and \textsc{small capitals} are useful for
names and terms that deserve a visual distinction without strong emphasis. Source code
and file names belong in \texttt{typewriter type}, and interface labels read best in
\textsf{sans serif}.

Font declarations act on the rest of the current group. The sentence {\bfseries this
phrase is set in bold by a declaration}, and {\itshape this one by an italic
declaration}, so that the surrounding text is unaffected. The older forms \bf and \it
work the same way inside braces, as in {\bf old-style bold} and {\it old-style italic}.

Accents are part of ordinary text. The word caf\'e takes an acute accent, the name
G\"odel takes an umlaut, and the phrase na\"\i ve reads as it should. Quotations are
typed as ``double quotes'' and the apostrophe in don't is a plain character, while
the dashes -- en and --- em -- have their own meanings.

Sentences that trail away end with \ldots{} and the typesetter spaces them as a unit.
The name of the system is written as \TeX{} in the running text and \LaTeX{} in the
document class that builds on it. Non-breaking spaces such as Figure~1 and Section~2
keep a label with its number, and a footnote gives a remark that would interrupt the
argument\footnote{A footnote is set at the foot of the page where its reference
appears, in a smaller size, and it is numbered in sequence with the other footnotes of
the document.}.

Explicit line breaks are useful in addresses and short verse:\\
the first line ends at the double backslash,\\
and the second line starts afresh.

Horizontal and vertical space can be inserted by hand. A paragraph that begins with
\noindent is not indented, and an explicit horizontal gap of \hspace{1em} separates
the two words in this sentence. Vertical space is more common between blocks, so the
next paragraph follows after \vspace{6pt} of extra room.

\noindent The paragraph above was started with the command \texttt{noindent}, so it
has no first-line indentation. The paragraphs that follow are indented as usual, which
is the default behaviour of the article class.

Ordinary paragraphs are separated by a blank line. The line breaking algorithm looks at
the whole paragraph before it commits to any break, which is why a change in one word
can move words several lines further down. The author rarely needs to intervene, but
the mechanisms exist, and the rest of this document uses them where they matter.

\subsection{Font declarations in context}\label{sec:fonts-2}

A declaration such as \textsf{bfseries} changes the font until the end of the enclosing
group. The next paragraph shows how groups limit the effect of such declarations, and
how the text returns to its original shape when the group closes.

{\bfseries Inside this group the text is bold, and a word such as \emph{this} still
changes to italic within bold text.} Outside the group the text returns to roman, as
the reader can see in this sentence. The same holds for {\itshape the slanted shape}
and for {\sffamily the sans serif family}, which all end where their braces end.

\subsection{Spacing and hyphenation}\label{sec:spacing-2}

Word spacing in a justified line is not constant. The stretch that a line receives is
shared among its interword glue, and the amount of stretch depends on how much room the
line has after its words are placed. A line with few words and a large gap is
noticeably loose, and a line with many words and little room may be tight. The
typesetter tries to keep all lines of a paragraph as even as possible, and it accepts
a slightly loose line when the alternative is a worse paragraph as a whole.

Long words that do not fit at the end of a line are hyphenated at points chosen from a
table of patterns. The table is built from word lists and the rules that govern where
a word may be divided, and it lets the typesetter break words such as typesetting,
hyphenation, algorithm and reproducibility at reasonable places. A hyphen at the end of
a line is a cost, and two such hyphens in a row are a larger cost, so the algorithm
avoids them where it can.

Some words must never be divided, and the author can say so by writing the word inside
a group that forbids breaks, or by joining the parts of a name with a tie. The name
Ludwig~van Beethoven keeps its two parts together, and a unit such as 12~pt does the
same for the number and the unit.

Sentences that end with a capital letter, such as the initials in J.~R.\ Smith, have a
different space after them, and the author must mark an end of sentence when the
capital is not the end of a sentence. The rule is simple, but it is easy to forget, and
the document below uses a backslash space in one case to show how it is written.

\section{Lists}\label{sec:lists-2}

Lists organise material that has no natural sentence structure. The three list
environments of the article class differ in how they mark their items and how they
number them. The first is the unordered list, which uses a bullet for each item.

\begin{itemize}
  \item The first item of a bulleted list is a plain sentence.
  \item The second item is longer and contains an \emph{emphasised} word, so that the
        wrapped line is indented to the same position as the text of the item.
  \item A third item may itself hold a sublist:
    \begin{itemize}
      \item the nested list uses a dash as its marker,
      \item and its items are indented further than the parent items.
    \end{itemize}
\end{itemize}

Numbered lists are used when the order of the items matters, for example in a
procedure. The following enumeration is nested two levels deep, and its first item
carries a label so that the text can refer to it later. The reference to that item
appears in Section~\ref{sec:tables-2}, which is written after this list.

\begin{enumerate}
  \item\label{item:first-2} Obtain the source files of the document and the packages that
        it loads.
  \item Compile the document once on the native path, which reads the auxiliary
        information of the previous run, if there is one.
        \begin{enumerate}
          \item The first sub-item describes a check that the author can make by hand.
          \item The second sub-item describes what happens when the check fails.
        \end{enumerate}
  \item Compare the result against the reference output and report any page that
        differs.
\end{enumerate}

The first step of the procedure is item~\ref{item:first-2} in the list above, and the
reader who wants the reasons behind the procedure will find them in
Section~\ref{sec:displays-2}.

Definition lists give a term followed by its explanation. The description environment
sets each term in bold and indents the explanation beneath it.

\begin{description}
  \item[Kerning] An adjustment of the space between two particular letters, such as
        the pair A and V, so that the shapes look evenly spaced.
  \item[Ligature] A single glyph that replaces a sequence of letters, such as the
        fi ligature in the word first.
  \item[Leading] The distance from one baseline to the next, measured in points.
\end{description}

The three kinds of list above are the core of most structured writing. An author who
needs a list inside a list simply nests the environment, and each level takes its own
label and its own indentation. The nesting depth in the article class is limited only
by the page, although a list nested more than three levels deep is usually a sign that
the material wants a table or a section instead.

Lists also interact with the page. A list is a sequence of paragraphs with extra space
between items, and the page builder may break the page between two items, as it may
between two paragraphs. The space at the top and bottom of a list is controlled by the
document class, so that a list opens and closes with the same visual distance as a
displayed equation.

\section{Displayed text}\label{sec:displays-2}

Some text is set apart from the body of the document to show that it is a quotation, a
centred title, or a block of literal material. Each environment has its own margins and
its own vertical spacing, and the environments below show the common cases.

\begin{quote}
A quotation environment sets its text with extra margin on both sides and without
paragraph indentation. It is intended for a short passage that the author wishes to
present as a separate block of prose, without the extra emphasis of a block quotation.
\end{quote}

\begin{quotation}
The quotation environment differs from the quote environment in that it indents the
first line of each of its paragraphs. A long passage that runs over several paragraphs
therefore keeps its internal structure, as a reader expects from a transcript of a
speech or a letter. This quotation has two paragraphs, so the second one can be seen to
start with an indent.

The second paragraph of the quotation shows the indentation of a continuing paragraph.
\end{quotation}

\begin{center}
Centred text on a line of its own\\
keeps each of its lines in the middle of the text block.
\end{center}

\begin{flushleft}
Flush left text keeps the left margin straight and lets the right margin
ragged, which suits material such as addresses and minutes of a meeting.
\end{flushleft}

\begin{flushright}
Flush right text keeps the right margin straight, which is seldom useful for
prose but is useful for a signature line.
\end{flushright}

Verbatim text is typeset exactly as it is typed, with spaces, line breaks and special
characters such as the percent sign and the braces kept as they stand:

\begin{verbatim}
  for i in 1..n {
      total += a[i] * b[i];   % dot product, not a comment
  }
\end{verbatim}

The following block uses a centring declaration inside a group, which keeps the
alignment local to the braces, and then a ragged right declaration on the next block.

{\centering A centred block made with a declaration, which ends at the closing
brace and leaves the surrounding text alone.\par}

{\raggedright A block whose lines are not stretched to the full measure, so the
right margin is ragged. The declaration ends with the group, and the next paragraph
returns to the usual justified setting of the text.\par}

Paragraphs of ordinary text after these blocks are justified again, as the reader can
see here.

\section{Tables}\label{sec:tables-2}

Tables set data in rows and columns. The tabular environment takes a column
specification, and each column is either left, centred or right aligned. Vertical bars
draw rules between columns. Section~\ref{sec:displays-2} explains the centring that surrounds each table in
this section.

\subsection{Basic alignment and rules}\label{sec:tables-basic-2}

\begin{center}
\begin{tabular}{l c r}
Method & Passes & Time (ms) \\
Greedy fill      & 1 & 0.4 \\
Knuth--Plass     & 2 & 3.1 \\
Optimal fit      & 3 & 4.7 \\
\end{tabular}
\end{center}

\begin{center}
\begin{tabular}{|l|c|r|}
Language & Typing & Speed \\
C        & static  & fast   \\
Python   & dynamic & slow   \\
Rust     & static  & fast   \\
\end{tabular}
\end{center}

\subsection{Padding and repeated columns}\label{sec:tables-advanced-2}

The column specification may fix the padding around a column with an empty at-sign
argument, which removes the space at one end of the table. Repeated columns are written
with a star argument, which expands to a run of identical columns.

\begin{center}
\begin{tabular}{@{}*{3}{c}@{}}
Left & Middle & Right \\
1 & 2 & 3 \\
4 & 5 & 6 \\
\end{tabular}
\end{center}

Each column keeps its own alignment in every row, so the numbers in the table below
line up on their right edges whatever the width of their labels.

\begin{center}
\begin{tabular}{l c r}
Run & Pages & Seconds \\
Alpha & 12 & 0.8 \\
Beta  & 29 & 1.9 \\
Gamma & 61 & 4.2 \\
\end{tabular}
\end{center}

Extra vertical space between rows is written after a row break with an optional length,
as in \texttt{[2pt]}. The table below uses that spacing to separate groups of rows.

\begin{center}
\begin{tabular}{l r}
Operation & Cost \\[2pt]
Read      & 3 \\
Write     & 7 \\[2pt]
Reference & 11 \\
\end{tabular}
\end{center}

Tables with many rows are common in reports, and the author who sets them by hand often
spends more time on alignment than on content. The tabular environment reduces that
work to a column specification and a list of rows. Each row ends with a double
backslash, and each cell is separated from its neighbour by an ampersand. The
alignment of each column is taken from the specification, so a change of one letter
moves the whole column without any other edit.

The tables so far show the alignment letters, the padding that the at-sign notation
controls and the repetition of a column pattern. Horizontal rules, partial rules and
cells that span several columns are set by the same machinery.

\subsection{Rules and spanning cells}\label{sec:tables-rules-2}
Table rules come in three kinds: a full rule across every column, a partial rule
across some of them, and a doubled rule made of two full rules in a row. A heading can
span all the columns of the table, and a single cell can span two of them.

\begin{center}
\begin{tabular}{|l|c|r|}
\hline
\multicolumn{3}{|c|}{\bfseries Measured costs} \\
\hline\hline
Phase & Calls & Time \\
\cline{2-3}
Scan    & 120 & 3.5 \\
Expand  & 940 & 12.0 \\
\multicolumn{2}{|l|}{Total, both phases} & 15.5 \\
\hline
\end{tabular}
\end{center}

\section{Floats}\label{sec:floats-2}

Floating material such as tables and figures is moved by the page builder to a position
where it fits. The placement specifier suggests where the float may go: \texttt{h} for
here, \texttt{t} for the top of a page, \texttt{b} for the bottom and \texttt{p} for a
page of floats. The author can give several letters at once, as in \texttt{[htbp]}, and
the page builder tries them in turn.

Table~\ref{tab:results-2} appears before the text that discusses it, so the reader meets
the reference first and the float later. The float is placed here by its own
specifier.

\begin{table}[h]
\centering
\begin{tabular}{l c r}
Input size & Lines & Pages \\
Short      & 40    & 1     \\
Medium     & 400   & 12    \\
Long       & 4000  & 120   \\
\end{tabular}
\caption{Pages produced for three input sizes.}\label{tab:results-2}
\end{table}

The results in Table~\ref{tab:results-2} grow with the input, and the page counts follow
the line counts closely. The next float is a figure with a caption long enough to wrap
onto a second line, and it is placed at the top of the page by its specifier.

\begin{figure}[t]
\centering
\mbox{\textsf{Input}\quad$\rightarrow$\quad\texttt{Parser}\quad$\rightarrow$\quad\textsf{Layout}\quad$\rightarrow$\quad\textsf{PDF}}
\caption{The stages of the native pipeline, from the characters of the source file to the
bytes of the output file, shown as a chain of boxes in the order in which the data flows
through them.}\label{fig:pipeline-2}
\end{figure}

Figure~\ref{fig:pipeline-2} shows the chain of stages that the native path runs in a
single pass. The figure caption is set as a paragraph because it runs over two lines.
A table placed at the bottom of a page is specified with \texttt{[b]}, as the next
example shows.

\begin{table}[b]
\centering
\begin{tabular}{c c}
Feature & State \\
Lists        & done \\
Tables       & done \\
\end{tabular}
\caption{Feature status at the time of writing.}
\end{table}

A float with both the top and the bottom permitted is written \texttt{[tb]}, and a float
that may go anywhere in the sequence is written \texttt{[htbp]}. The next figure uses
the latter specifier and a label that the text refers to from a later section.

\begin{figure}[htbp]
\centering
\mbox{\texttt{figure}} \quad \mbox{Text placed in a box}
\caption{A short caption.}\label{fig:box-2}
\end{figure}

Figure~\ref{fig:box-2} is referred to here, after the float itself. The references to the
tables and figures of this section are resolved in the second pass, which is the reason
that the reader sees the correct numbers even for references that come before their
floats.

A float is a block of material that the page builder may move away from the place where
it appears in the source. Its purpose is to keep large items, such as a table of
results or a diagram, from splitting across a page break or from leaving a large gap
at the foot of a page. The page builder reads the placement specifier of the float and
considers each candidate position in turn. If no candidate fits, the float waits for
the next page, and the material that follows it continues to fill the current page.

A caption belongs to its float and is numbered with it. The number of a table and the
number of a figure are kept in separate counters, so that a document may contain a
table 1 and a figure 1 without conflict. The label placed after the caption records the
number, and the reference in the text picks up the number from the label when the
document is compiled. This is the reason that a reference may come before its float in
the source, as the references to the tables above do.

\section{Mathematics}\label{sec:math-2}

Mathematical notation is set in math mode, which uses its own fonts and spacing rules.
Inline formulae are written between dollar signs. The quantity $\frac{a+b}{c-d}$ is a
fraction, and $\sqrt{x^2+1}$ is a square root. A superscript such as $x^{2}$ and a
subscript such as $a_{i}$ are common, and both can appear together, as in $x_{i}^{2}$.

Greek letters are typed by name, so $\alpha$, $\beta$ and $\gamma$ are available, and so
is the capital $\Gamma$. The sum $\sum_{k=1}^{n} k$ and the integral $\int_0^1 x\,dx$
appear inline with their limits set beside them, as is customary in running text. The
limit $\lim_{n\to\infty} \frac{1}{n} = 0$ also appears inline.

Sets are typeset with blackboard and calligraphic letters. The set of real numbers is
$\mathbb{R}$, and a function that bounds another is written $f(n) \in \mathcal{O}(n^2)$.
A vector is written in bold, as $\mathbf{v}$, and its length is $\|\mathbf{v}\|$. The
parentheses in $\left( \frac{x}{y} \right)$ grow with the fraction they enclose.

Displayed formulae are set on their own line, centred, with a larger style than inline
formulae. The formula below gives the sum of the first $n$ integers.
\[
  \sum_{k=1}^{n} k = \frac{n(n+1)}{2}
\]

Numbered equations carry a label, so that the text can refer to them by number. The
first equation of this section is the definition of the Gaussian integral, and the
second is the quadratic formula, which is used in Section~\ref{sec:lists-2}.

\begin{equation}\label{eq:gauss-2}
  \int_{-\infty}^{\infty} e^{-x^{2}}\,dx = \sqrt{\pi}
\end{equation}

\begin{equation}\label{eq:quadratic-2}
  x = \frac{-b \pm \sqrt{b^{2}-4ac}}{2a}
\end{equation}

Equation~\ref{eq:gauss-2} gives the Gaussian integral, and
Equation~\ref{eq:quadratic-2} gives the roots of a quadratic polynomial. The coefficients
$a$, $b$ and $c$ in Equation~\ref{eq:quadratic-2} are real numbers with $a \neq 0$.

A last display shows a sum over a product, with a subscript that is itself a formula:
\[
  \prod_{i=1}^{m} \left( 1 + \frac{1}{i^{2}} \right) \leq e^{\pi^{2}/6}
\]

Mathematical formulae are read as a sequence of symbols with a structure, and the layout
rules give each symbol a size and a position. Inline formulae are set in the text style,
which keeps the lines of a paragraph at a uniform height, while displayed formulae are
set in the display style, which allows larger operators and fractions with larger
numerators and denominators. The reader sees the difference in the sums and products
above, where the limits sit above and below the symbol in the display, and beside it in
the running text.

The quadratic formula in Equation~\ref{eq:quadratic-2} is a familiar example of a
fraction with a root in its numerator. The square root of the discriminant
$b^{2}-4ac$ is set with a radical sign that grows with its contents, and the
fraction line extends beneath the whole numerator. The parentheses of the next example
are sized in the same way, so that the reader can see the structure at a glance.

\subsection{Displayed derivations}\label{sec:derivations-2}

A derivation shows a chain of equalities, each following from the one before. Each line
of the chain is written as a separate display, and the text between the displays gives
the reason for each step. The first derivation below expands the square of a binomial,
and the second uses the result to evaluate a finite sum of squares.
\[
  (a+b)^{2} = a^{2} + 2ab + b^{2}
\]

Applying the same expansion to the terms of a sum gives a closed form for the sum of the
squares of the first $n$ integers, which is the second formula of this subsection:
\[
  \sum_{k=1}^{n} k^{2} = \frac{n(n+1)(2n+1)}{6}
\]

The formula holds for every positive integer $n$, and it can be proved by induction on
$n$, as the reader may check for the first few values. Each step of the induction
uses the identity above and the fact that the sum for $n+1$ is the sum for $n$ with one
more term added.

The third display is a limit that defines the constant $e$. It is set on its own line
because it is a result that the text refers to, and the reader may wish to compare it
with the bound given in the last display of the previous section:
\[
  e = \lim_{n\to\infty} \left( 1 + \frac{1}{n} \right)^{n}
\]

Derivations that are long are best broken into several displays, each short enough to
fit on one line of the text block, so that the reader can follow the argument without
turning the page.

\section{Bibliography and citations}\label{sec:refs-2}

The native path follows the algorithms of the typesetting system closely, and the
descriptions of those algorithms are found in the standard references. The line
breaking method used by the paragraph builder is described by Knuth and
Plass~\cite{knuth81-2}, while the page builder and the float placement are described in
the reference manual~\cite{knuth86-2}. The layout of mathematical formulae follows the
rules of the appendix on mathematics in that same manual. The chain of stages in
Figure~\ref{fig:pipeline-2} is a common design for compilers in general, as Aho and
colleagues describe~\cite{aho06-2}.

\begin{thebibliography}{9}
\bibitem{knuth81-2}
  Donald E.~Knuth and Michael F.~Plass.
  \emph{Breaking paragraphs into lines}.
  Software: Practice and Experience, 11(11), 1981.

\bibitem{knuth86-2}
  Donald E.~Knuth.
  \emph{The \TeX book}.
  Addison-Wesley, 1986.

\bibitem{aho06-2}
  Alfred V.~Aho, Monica S.~Lam, Ravi Sethi and Jeffrey D.~Ullman.
  \emph{Compilers: Principles, Techniques, and Tools}.
  Addison-Wesley, second edition, 2006.
\end{thebibliography}



\section{Running text}\label{sec:text-3}

Typesetting is the business of turning a stream of characters into a page. A document
author writes paragraphs of plain prose, and the typesetter decides where lines break,
how much space separates words, and where the page ends. The quality of the result
depends on many small decisions, most of which the author never sees. This section
collects the inline devices that an author reaches for while writing prose: emphasis,
font families, font shapes, accents and punctuation.

Emphasis is usually marked with \emph{italic} text, which in an upright paragraph
reverses to roman when nested inside an italic context. Strong statements can be set in
\textbf{bold}, while \textit{explicit italic} and \textsc{small capitals} are useful for
names and terms that deserve a visual distinction without strong emphasis. Source code
and file names belong in \texttt{typewriter type}, and interface labels read best in
\textsf{sans serif}.

Font declarations act on the rest of the current group. The sentence {\bfseries this
phrase is set in bold by a declaration}, and {\itshape this one by an italic
declaration}, so that the surrounding text is unaffected. The older forms \bf and \it
work the same way inside braces, as in {\bf old-style bold} and {\it old-style italic}.

Accents are part of ordinary text. The word caf\'e takes an acute accent, the name
G\"odel takes an umlaut, and the phrase na\"\i ve reads as it should. Quotations are
typed as ``double quotes'' and the apostrophe in don't is a plain character, while
the dashes -- en and --- em -- have their own meanings.

Sentences that trail away end with \ldots{} and the typesetter spaces them as a unit.
The name of the system is written as \TeX{} in the running text and \LaTeX{} in the
document class that builds on it. Non-breaking spaces such as Figure~1 and Section~2
keep a label with its number, and a footnote gives a remark that would interrupt the
argument\footnote{A footnote is set at the foot of the page where its reference
appears, in a smaller size, and it is numbered in sequence with the other footnotes of
the document.}.

Explicit line breaks are useful in addresses and short verse:\\
the first line ends at the double backslash,\\
and the second line starts afresh.

Horizontal and vertical space can be inserted by hand. A paragraph that begins with
\noindent is not indented, and an explicit horizontal gap of \hspace{1em} separates
the two words in this sentence. Vertical space is more common between blocks, so the
next paragraph follows after \vspace{6pt} of extra room.

\noindent The paragraph above was started with the command \texttt{noindent}, so it
has no first-line indentation. The paragraphs that follow are indented as usual, which
is the default behaviour of the article class.

Ordinary paragraphs are separated by a blank line. The line breaking algorithm looks at
the whole paragraph before it commits to any break, which is why a change in one word
can move words several lines further down. The author rarely needs to intervene, but
the mechanisms exist, and the rest of this document uses them where they matter.

\subsection{Font declarations in context}\label{sec:fonts-3}

A declaration such as \textsf{bfseries} changes the font until the end of the enclosing
group. The next paragraph shows how groups limit the effect of such declarations, and
how the text returns to its original shape when the group closes.

{\bfseries Inside this group the text is bold, and a word such as \emph{this} still
changes to italic within bold text.} Outside the group the text returns to roman, as
the reader can see in this sentence. The same holds for {\itshape the slanted shape}
and for {\sffamily the sans serif family}, which all end where their braces end.

\subsection{Spacing and hyphenation}\label{sec:spacing-3}

Word spacing in a justified line is not constant. The stretch that a line receives is
shared among its interword glue, and the amount of stretch depends on how much room the
line has after its words are placed. A line with few words and a large gap is
noticeably loose, and a line with many words and little room may be tight. The
typesetter tries to keep all lines of a paragraph as even as possible, and it accepts
a slightly loose line when the alternative is a worse paragraph as a whole.

Long words that do not fit at the end of a line are hyphenated at points chosen from a
table of patterns. The table is built from word lists and the rules that govern where
a word may be divided, and it lets the typesetter break words such as typesetting,
hyphenation, algorithm and reproducibility at reasonable places. A hyphen at the end of
a line is a cost, and two such hyphens in a row are a larger cost, so the algorithm
avoids them where it can.

Some words must never be divided, and the author can say so by writing the word inside
a group that forbids breaks, or by joining the parts of a name with a tie. The name
Ludwig~van Beethoven keeps its two parts together, and a unit such as 12~pt does the
same for the number and the unit.

Sentences that end with a capital letter, such as the initials in J.~R.\ Smith, have a
different space after them, and the author must mark an end of sentence when the
capital is not the end of a sentence. The rule is simple, but it is easy to forget, and
the document below uses a backslash space in one case to show how it is written.

\section{Lists}\label{sec:lists-3}

Lists organise material that has no natural sentence structure. The three list
environments of the article class differ in how they mark their items and how they
number them. The first is the unordered list, which uses a bullet for each item.

\begin{itemize}
  \item The first item of a bulleted list is a plain sentence.
  \item The second item is longer and contains an \emph{emphasised} word, so that the
        wrapped line is indented to the same position as the text of the item.
  \item A third item may itself hold a sublist:
    \begin{itemize}
      \item the nested list uses a dash as its marker,
      \item and its items are indented further than the parent items.
    \end{itemize}
\end{itemize}

Numbered lists are used when the order of the items matters, for example in a
procedure. The following enumeration is nested two levels deep, and its first item
carries a label so that the text can refer to it later. The reference to that item
appears in Section~\ref{sec:tables-3}, which is written after this list.

\begin{enumerate}
  \item\label{item:first-3} Obtain the source files of the document and the packages that
        it loads.
  \item Compile the document once on the native path, which reads the auxiliary
        information of the previous run, if there is one.
        \begin{enumerate}
          \item The first sub-item describes a check that the author can make by hand.
          \item The second sub-item describes what happens when the check fails.
        \end{enumerate}
  \item Compare the result against the reference output and report any page that
        differs.
\end{enumerate}

The first step of the procedure is item~\ref{item:first-3} in the list above, and the
reader who wants the reasons behind the procedure will find them in
Section~\ref{sec:displays-3}.

Definition lists give a term followed by its explanation. The description environment
sets each term in bold and indents the explanation beneath it.

\begin{description}
  \item[Kerning] An adjustment of the space between two particular letters, such as
        the pair A and V, so that the shapes look evenly spaced.
  \item[Ligature] A single glyph that replaces a sequence of letters, such as the
        fi ligature in the word first.
  \item[Leading] The distance from one baseline to the next, measured in points.
\end{description}

The three kinds of list above are the core of most structured writing. An author who
needs a list inside a list simply nests the environment, and each level takes its own
label and its own indentation. The nesting depth in the article class is limited only
by the page, although a list nested more than three levels deep is usually a sign that
the material wants a table or a section instead.

Lists also interact with the page. A list is a sequence of paragraphs with extra space
between items, and the page builder may break the page between two items, as it may
between two paragraphs. The space at the top and bottom of a list is controlled by the
document class, so that a list opens and closes with the same visual distance as a
displayed equation.

\section{Displayed text}\label{sec:displays-3}

Some text is set apart from the body of the document to show that it is a quotation, a
centred title, or a block of literal material. Each environment has its own margins and
its own vertical spacing, and the environments below show the common cases.

\begin{quote}
A quotation environment sets its text with extra margin on both sides and without
paragraph indentation. It is intended for a short passage that the author wishes to
present as a separate block of prose, without the extra emphasis of a block quotation.
\end{quote}

\begin{quotation}
The quotation environment differs from the quote environment in that it indents the
first line of each of its paragraphs. A long passage that runs over several paragraphs
therefore keeps its internal structure, as a reader expects from a transcript of a
speech or a letter. This quotation has two paragraphs, so the second one can be seen to
start with an indent.

The second paragraph of the quotation shows the indentation of a continuing paragraph.
\end{quotation}

\begin{center}
Centred text on a line of its own\\
keeps each of its lines in the middle of the text block.
\end{center}

\begin{flushleft}
Flush left text keeps the left margin straight and lets the right margin
ragged, which suits material such as addresses and minutes of a meeting.
\end{flushleft}

\begin{flushright}
Flush right text keeps the right margin straight, which is seldom useful for
prose but is useful for a signature line.
\end{flushright}

Verbatim text is typeset exactly as it is typed, with spaces, line breaks and special
characters such as the percent sign and the braces kept as they stand:

\begin{verbatim}
  for i in 1..n {
      total += a[i] * b[i];   % dot product, not a comment
  }
\end{verbatim}

The following block uses a centring declaration inside a group, which keeps the
alignment local to the braces, and then a ragged right declaration on the next block.

{\centering A centred block made with a declaration, which ends at the closing
brace and leaves the surrounding text alone.\par}

{\raggedright A block whose lines are not stretched to the full measure, so the
right margin is ragged. The declaration ends with the group, and the next paragraph
returns to the usual justified setting of the text.\par}

Paragraphs of ordinary text after these blocks are justified again, as the reader can
see here.

\section{Tables}\label{sec:tables-3}

Tables set data in rows and columns. The tabular environment takes a column
specification, and each column is either left, centred or right aligned. Vertical bars
draw rules between columns. Section~\ref{sec:displays-3} explains the centring that surrounds each table in
this section.

\subsection{Basic alignment and rules}\label{sec:tables-basic-3}

\begin{center}
\begin{tabular}{l c r}
Method & Passes & Time (ms) \\
Greedy fill      & 1 & 0.4 \\
Knuth--Plass     & 2 & 3.1 \\
Optimal fit      & 3 & 4.7 \\
\end{tabular}
\end{center}

\begin{center}
\begin{tabular}{|l|c|r|}
Language & Typing & Speed \\
C        & static  & fast   \\
Python   & dynamic & slow   \\
Rust     & static  & fast   \\
\end{tabular}
\end{center}

\subsection{Padding and repeated columns}\label{sec:tables-advanced-3}

The column specification may fix the padding around a column with an empty at-sign
argument, which removes the space at one end of the table. Repeated columns are written
with a star argument, which expands to a run of identical columns.

\begin{center}
\begin{tabular}{@{}*{3}{c}@{}}
Left & Middle & Right \\
1 & 2 & 3 \\
4 & 5 & 6 \\
\end{tabular}
\end{center}

Each column keeps its own alignment in every row, so the numbers in the table below
line up on their right edges whatever the width of their labels.

\begin{center}
\begin{tabular}{l c r}
Run & Pages & Seconds \\
Alpha & 12 & 0.8 \\
Beta  & 29 & 1.9 \\
Gamma & 61 & 4.2 \\
\end{tabular}
\end{center}

Extra vertical space between rows is written after a row break with an optional length,
as in \texttt{[2pt]}. The table below uses that spacing to separate groups of rows.

\begin{center}
\begin{tabular}{l r}
Operation & Cost \\[2pt]
Read      & 3 \\
Write     & 7 \\[2pt]
Reference & 11 \\
\end{tabular}
\end{center}

Tables with many rows are common in reports, and the author who sets them by hand often
spends more time on alignment than on content. The tabular environment reduces that
work to a column specification and a list of rows. Each row ends with a double
backslash, and each cell is separated from its neighbour by an ampersand. The
alignment of each column is taken from the specification, so a change of one letter
moves the whole column without any other edit.

The tables so far show the alignment letters, the padding that the at-sign notation
controls and the repetition of a column pattern. Horizontal rules, partial rules and
cells that span several columns are set by the same machinery.

\subsection{Rules and spanning cells}\label{sec:tables-rules-3}
Table rules come in three kinds: a full rule across every column, a partial rule
across some of them, and a doubled rule made of two full rules in a row. A heading can
span all the columns of the table, and a single cell can span two of them.

\begin{center}
\begin{tabular}{|l|c|r|}
\hline
\multicolumn{3}{|c|}{\bfseries Measured costs} \\
\hline\hline
Phase & Calls & Time \\
\cline{2-3}
Scan    & 120 & 3.5 \\
Expand  & 940 & 12.0 \\
\multicolumn{2}{|l|}{Total, both phases} & 15.5 \\
\hline
\end{tabular}
\end{center}

\section{Floats}\label{sec:floats-3}

Floating material such as tables and figures is moved by the page builder to a position
where it fits. The placement specifier suggests where the float may go: \texttt{h} for
here, \texttt{t} for the top of a page, \texttt{b} for the bottom and \texttt{p} for a
page of floats. The author can give several letters at once, as in \texttt{[htbp]}, and
the page builder tries them in turn.

Table~\ref{tab:results-3} appears before the text that discusses it, so the reader meets
the reference first and the float later. The float is placed here by its own
specifier.

\begin{table}[h]
\centering
\begin{tabular}{l c r}
Input size & Lines & Pages \\
Short      & 40    & 1     \\
Medium     & 400   & 12    \\
Long       & 4000  & 120   \\
\end{tabular}
\caption{Pages produced for three input sizes.}\label{tab:results-3}
\end{table}

The results in Table~\ref{tab:results-3} grow with the input, and the page counts follow
the line counts closely. The next float is a figure with a caption long enough to wrap
onto a second line, and it is placed at the top of the page by its specifier.

\begin{figure}[t]
\centering
\mbox{\textsf{Input}\quad$\rightarrow$\quad\texttt{Parser}\quad$\rightarrow$\quad\textsf{Layout}\quad$\rightarrow$\quad\textsf{PDF}}
\caption{The stages of the native pipeline, from the characters of the source file to the
bytes of the output file, shown as a chain of boxes in the order in which the data flows
through them.}\label{fig:pipeline-3}
\end{figure}

Figure~\ref{fig:pipeline-3} shows the chain of stages that the native path runs in a
single pass. The figure caption is set as a paragraph because it runs over two lines.
A table placed at the bottom of a page is specified with \texttt{[b]}, as the next
example shows.

\begin{table}[b]
\centering
\begin{tabular}{c c}
Feature & State \\
Lists        & done \\
Tables       & done \\
\end{tabular}
\caption{Feature status at the time of writing.}
\end{table}

A float with both the top and the bottom permitted is written \texttt{[tb]}, and a float
that may go anywhere in the sequence is written \texttt{[htbp]}. The next figure uses
the latter specifier and a label that the text refers to from a later section.

\begin{figure}[htbp]
\centering
\mbox{\texttt{figure}} \quad \mbox{Text placed in a box}
\caption{A short caption.}\label{fig:box-3}
\end{figure}

Figure~\ref{fig:box-3} is referred to here, after the float itself. The references to the
tables and figures of this section are resolved in the second pass, which is the reason
that the reader sees the correct numbers even for references that come before their
floats.

A float is a block of material that the page builder may move away from the place where
it appears in the source. Its purpose is to keep large items, such as a table of
results or a diagram, from splitting across a page break or from leaving a large gap
at the foot of a page. The page builder reads the placement specifier of the float and
considers each candidate position in turn. If no candidate fits, the float waits for
the next page, and the material that follows it continues to fill the current page.

A caption belongs to its float and is numbered with it. The number of a table and the
number of a figure are kept in separate counters, so that a document may contain a
table 1 and a figure 1 without conflict. The label placed after the caption records the
number, and the reference in the text picks up the number from the label when the
document is compiled. This is the reason that a reference may come before its float in
the source, as the references to the tables above do.

\section{Mathematics}\label{sec:math-3}

Mathematical notation is set in math mode, which uses its own fonts and spacing rules.
Inline formulae are written between dollar signs. The quantity $\frac{a+b}{c-d}$ is a
fraction, and $\sqrt{x^2+1}$ is a square root. A superscript such as $x^{2}$ and a
subscript such as $a_{i}$ are common, and both can appear together, as in $x_{i}^{2}$.

Greek letters are typed by name, so $\alpha$, $\beta$ and $\gamma$ are available, and so
is the capital $\Gamma$. The sum $\sum_{k=1}^{n} k$ and the integral $\int_0^1 x\,dx$
appear inline with their limits set beside them, as is customary in running text. The
limit $\lim_{n\to\infty} \frac{1}{n} = 0$ also appears inline.

Sets are typeset with blackboard and calligraphic letters. The set of real numbers is
$\mathbb{R}$, and a function that bounds another is written $f(n) \in \mathcal{O}(n^2)$.
A vector is written in bold, as $\mathbf{v}$, and its length is $\|\mathbf{v}\|$. The
parentheses in $\left( \frac{x}{y} \right)$ grow with the fraction they enclose.

Displayed formulae are set on their own line, centred, with a larger style than inline
formulae. The formula below gives the sum of the first $n$ integers.
\[
  \sum_{k=1}^{n} k = \frac{n(n+1)}{2}
\]

Numbered equations carry a label, so that the text can refer to them by number. The
first equation of this section is the definition of the Gaussian integral, and the
second is the quadratic formula, which is used in Section~\ref{sec:lists-3}.

\begin{equation}\label{eq:gauss-3}
  \int_{-\infty}^{\infty} e^{-x^{2}}\,dx = \sqrt{\pi}
\end{equation}

\begin{equation}\label{eq:quadratic-3}
  x = \frac{-b \pm \sqrt{b^{2}-4ac}}{2a}
\end{equation}

Equation~\ref{eq:gauss-3} gives the Gaussian integral, and
Equation~\ref{eq:quadratic-3} gives the roots of a quadratic polynomial. The coefficients
$a$, $b$ and $c$ in Equation~\ref{eq:quadratic-3} are real numbers with $a \neq 0$.

A last display shows a sum over a product, with a subscript that is itself a formula:
\[
  \prod_{i=1}^{m} \left( 1 + \frac{1}{i^{2}} \right) \leq e^{\pi^{2}/6}
\]

Mathematical formulae are read as a sequence of symbols with a structure, and the layout
rules give each symbol a size and a position. Inline formulae are set in the text style,
which keeps the lines of a paragraph at a uniform height, while displayed formulae are
set in the display style, which allows larger operators and fractions with larger
numerators and denominators. The reader sees the difference in the sums and products
above, where the limits sit above and below the symbol in the display, and beside it in
the running text.

The quadratic formula in Equation~\ref{eq:quadratic-3} is a familiar example of a
fraction with a root in its numerator. The square root of the discriminant
$b^{2}-4ac$ is set with a radical sign that grows with its contents, and the
fraction line extends beneath the whole numerator. The parentheses of the next example
are sized in the same way, so that the reader can see the structure at a glance.

\subsection{Displayed derivations}\label{sec:derivations-3}

A derivation shows a chain of equalities, each following from the one before. Each line
of the chain is written as a separate display, and the text between the displays gives
the reason for each step. The first derivation below expands the square of a binomial,
and the second uses the result to evaluate a finite sum of squares.
\[
  (a+b)^{2} = a^{2} + 2ab + b^{2}
\]

Applying the same expansion to the terms of a sum gives a closed form for the sum of the
squares of the first $n$ integers, which is the second formula of this subsection:
\[
  \sum_{k=1}^{n} k^{2} = \frac{n(n+1)(2n+1)}{6}
\]

The formula holds for every positive integer $n$, and it can be proved by induction on
$n$, as the reader may check for the first few values. Each step of the induction
uses the identity above and the fact that the sum for $n+1$ is the sum for $n$ with one
more term added.

The third display is a limit that defines the constant $e$. It is set on its own line
because it is a result that the text refers to, and the reader may wish to compare it
with the bound given in the last display of the previous section:
\[
  e = \lim_{n\to\infty} \left( 1 + \frac{1}{n} \right)^{n}
\]

Derivations that are long are best broken into several displays, each short enough to
fit on one line of the text block, so that the reader can follow the argument without
turning the page.

\section{Bibliography and citations}\label{sec:refs-3}

The native path follows the algorithms of the typesetting system closely, and the
descriptions of those algorithms are found in the standard references. The line
breaking method used by the paragraph builder is described by Knuth and
Plass~\cite{knuth81-3}, while the page builder and the float placement are described in
the reference manual~\cite{knuth86-3}. The layout of mathematical formulae follows the
rules of the appendix on mathematics in that same manual. The chain of stages in
Figure~\ref{fig:pipeline-3} is a common design for compilers in general, as Aho and
colleagues describe~\cite{aho06-3}.

\begin{thebibliography}{9}
\bibitem{knuth81-3}
  Donald E.~Knuth and Michael F.~Plass.
  \emph{Breaking paragraphs into lines}.
  Software: Practice and Experience, 11(11), 1981.

\bibitem{knuth86-3}
  Donald E.~Knuth.
  \emph{The \TeX book}.
  Addison-Wesley, 1986.

\bibitem{aho06-3}
  Alfred V.~Aho, Monica S.~Lam, Ravi Sethi and Jeffrey D.~Ullman.
  \emph{Compilers: Principles, Techniques, and Tools}.
  Addison-Wesley, second edition, 2006.
\end{thebibliography}



\section{Running text}\label{sec:text-4}

Typesetting is the business of turning a stream of characters into a page. A document
author writes paragraphs of plain prose, and the typesetter decides where lines break,
how much space separates words, and where the page ends. The quality of the result
depends on many small decisions, most of which the author never sees. This section
collects the inline devices that an author reaches for while writing prose: emphasis,
font families, font shapes, accents and punctuation.

Emphasis is usually marked with \emph{italic} text, which in an upright paragraph
reverses to roman when nested inside an italic context. Strong statements can be set in
\textbf{bold}, while \textit{explicit italic} and \textsc{small capitals} are useful for
names and terms that deserve a visual distinction without strong emphasis. Source code
and file names belong in \texttt{typewriter type}, and interface labels read best in
\textsf{sans serif}.

Font declarations act on the rest of the current group. The sentence {\bfseries this
phrase is set in bold by a declaration}, and {\itshape this one by an italic
declaration}, so that the surrounding text is unaffected. The older forms \bf and \it
work the same way inside braces, as in {\bf old-style bold} and {\it old-style italic}.

Accents are part of ordinary text. The word caf\'e takes an acute accent, the name
G\"odel takes an umlaut, and the phrase na\"\i ve reads as it should. Quotations are
typed as ``double quotes'' and the apostrophe in don't is a plain character, while
the dashes -- en and --- em -- have their own meanings.

Sentences that trail away end with \ldots{} and the typesetter spaces them as a unit.
The name of the system is written as \TeX{} in the running text and \LaTeX{} in the
document class that builds on it. Non-breaking spaces such as Figure~1 and Section~2
keep a label with its number, and a footnote gives a remark that would interrupt the
argument\footnote{A footnote is set at the foot of the page where its reference
appears, in a smaller size, and it is numbered in sequence with the other footnotes of
the document.}.

Explicit line breaks are useful in addresses and short verse:\\
the first line ends at the double backslash,\\
and the second line starts afresh.

Horizontal and vertical space can be inserted by hand. A paragraph that begins with
\noindent is not indented, and an explicit horizontal gap of \hspace{1em} separates
the two words in this sentence. Vertical space is more common between blocks, so the
next paragraph follows after \vspace{6pt} of extra room.

\noindent The paragraph above was started with the command \texttt{noindent}, so it
has no first-line indentation. The paragraphs that follow are indented as usual, which
is the default behaviour of the article class.

Ordinary paragraphs are separated by a blank line. The line breaking algorithm looks at
the whole paragraph before it commits to any break, which is why a change in one word
can move words several lines further down. The author rarely needs to intervene, but
the mechanisms exist, and the rest of this document uses them where they matter.

\subsection{Font declarations in context}\label{sec:fonts-4}

A declaration such as \textsf{bfseries} changes the font until the end of the enclosing
group. The next paragraph shows how groups limit the effect of such declarations, and
how the text returns to its original shape when the group closes.

{\bfseries Inside this group the text is bold, and a word such as \emph{this} still
changes to italic within bold text.} Outside the group the text returns to roman, as
the reader can see in this sentence. The same holds for {\itshape the slanted shape}
and for {\sffamily the sans serif family}, which all end where their braces end.

\subsection{Spacing and hyphenation}\label{sec:spacing-4}

Word spacing in a justified line is not constant. The stretch that a line receives is
shared among its interword glue, and the amount of stretch depends on how much room the
line has after its words are placed. A line with few words and a large gap is
noticeably loose, and a line with many words and little room may be tight. The
typesetter tries to keep all lines of a paragraph as even as possible, and it accepts
a slightly loose line when the alternative is a worse paragraph as a whole.

Long words that do not fit at the end of a line are hyphenated at points chosen from a
table of patterns. The table is built from word lists and the rules that govern where
a word may be divided, and it lets the typesetter break words such as typesetting,
hyphenation, algorithm and reproducibility at reasonable places. A hyphen at the end of
a line is a cost, and two such hyphens in a row are a larger cost, so the algorithm
avoids them where it can.

Some words must never be divided, and the author can say so by writing the word inside
a group that forbids breaks, or by joining the parts of a name with a tie. The name
Ludwig~van Beethoven keeps its two parts together, and a unit such as 12~pt does the
same for the number and the unit.

Sentences that end with a capital letter, such as the initials in J.~R.\ Smith, have a
different space after them, and the author must mark an end of sentence when the
capital is not the end of a sentence. The rule is simple, but it is easy to forget, and
the document below uses a backslash space in one case to show how it is written.

\section{Lists}\label{sec:lists-4}

Lists organise material that has no natural sentence structure. The three list
environments of the article class differ in how they mark their items and how they
number them. The first is the unordered list, which uses a bullet for each item.

\begin{itemize}
  \item The first item of a bulleted list is a plain sentence.
  \item The second item is longer and contains an \emph{emphasised} word, so that the
        wrapped line is indented to the same position as the text of the item.
  \item A third item may itself hold a sublist:
    \begin{itemize}
      \item the nested list uses a dash as its marker,
      \item and its items are indented further than the parent items.
    \end{itemize}
\end{itemize}

Numbered lists are used when the order of the items matters, for example in a
procedure. The following enumeration is nested two levels deep, and its first item
carries a label so that the text can refer to it later. The reference to that item
appears in Section~\ref{sec:tables-4}, which is written after this list.

\begin{enumerate}
  \item\label{item:first-4} Obtain the source files of the document and the packages that
        it loads.
  \item Compile the document once on the native path, which reads the auxiliary
        information of the previous run, if there is one.
        \begin{enumerate}
          \item The first sub-item describes a check that the author can make by hand.
          \item The second sub-item describes what happens when the check fails.
        \end{enumerate}
  \item Compare the result against the reference output and report any page that
        differs.
\end{enumerate}

The first step of the procedure is item~\ref{item:first-4} in the list above, and the
reader who wants the reasons behind the procedure will find them in
Section~\ref{sec:displays-4}.

Definition lists give a term followed by its explanation. The description environment
sets each term in bold and indents the explanation beneath it.

\begin{description}
  \item[Kerning] An adjustment of the space between two particular letters, such as
        the pair A and V, so that the shapes look evenly spaced.
  \item[Ligature] A single glyph that replaces a sequence of letters, such as the
        fi ligature in the word first.
  \item[Leading] The distance from one baseline to the next, measured in points.
\end{description}

The three kinds of list above are the core of most structured writing. An author who
needs a list inside a list simply nests the environment, and each level takes its own
label and its own indentation. The nesting depth in the article class is limited only
by the page, although a list nested more than three levels deep is usually a sign that
the material wants a table or a section instead.

Lists also interact with the page. A list is a sequence of paragraphs with extra space
between items, and the page builder may break the page between two items, as it may
between two paragraphs. The space at the top and bottom of a list is controlled by the
document class, so that a list opens and closes with the same visual distance as a
displayed equation.

\section{Displayed text}\label{sec:displays-4}

Some text is set apart from the body of the document to show that it is a quotation, a
centred title, or a block of literal material. Each environment has its own margins and
its own vertical spacing, and the environments below show the common cases.

\begin{quote}
A quotation environment sets its text with extra margin on both sides and without
paragraph indentation. It is intended for a short passage that the author wishes to
present as a separate block of prose, without the extra emphasis of a block quotation.
\end{quote}

\begin{quotation}
The quotation environment differs from the quote environment in that it indents the
first line of each of its paragraphs. A long passage that runs over several paragraphs
therefore keeps its internal structure, as a reader expects from a transcript of a
speech or a letter. This quotation has two paragraphs, so the second one can be seen to
start with an indent.

The second paragraph of the quotation shows the indentation of a continuing paragraph.
\end{quotation}

\begin{center}
Centred text on a line of its own\\
keeps each of its lines in the middle of the text block.
\end{center}

\begin{flushleft}
Flush left text keeps the left margin straight and lets the right margin
ragged, which suits material such as addresses and minutes of a meeting.
\end{flushleft}

\begin{flushright}
Flush right text keeps the right margin straight, which is seldom useful for
prose but is useful for a signature line.
\end{flushright}

Verbatim text is typeset exactly as it is typed, with spaces, line breaks and special
characters such as the percent sign and the braces kept as they stand:

\begin{verbatim}
  for i in 1..n {
      total += a[i] * b[i];   % dot product, not a comment
  }
\end{verbatim}

The following block uses a centring declaration inside a group, which keeps the
alignment local to the braces, and then a ragged right declaration on the next block.

{\centering A centred block made with a declaration, which ends at the closing
brace and leaves the surrounding text alone.\par}

{\raggedright A block whose lines are not stretched to the full measure, so the
right margin is ragged. The declaration ends with the group, and the next paragraph
returns to the usual justified setting of the text.\par}

Paragraphs of ordinary text after these blocks are justified again, as the reader can
see here.

\section{Tables}\label{sec:tables-4}

Tables set data in rows and columns. The tabular environment takes a column
specification, and each column is either left, centred or right aligned. Vertical bars
draw rules between columns. Section~\ref{sec:displays-4} explains the centring that surrounds each table in
this section.

\subsection{Basic alignment and rules}\label{sec:tables-basic-4}

\begin{center}
\begin{tabular}{l c r}
Method & Passes & Time (ms) \\
Greedy fill      & 1 & 0.4 \\
Knuth--Plass     & 2 & 3.1 \\
Optimal fit      & 3 & 4.7 \\
\end{tabular}
\end{center}

\begin{center}
\begin{tabular}{|l|c|r|}
Language & Typing & Speed \\
C        & static  & fast   \\
Python   & dynamic & slow   \\
Rust     & static  & fast   \\
\end{tabular}
\end{center}

\subsection{Padding and repeated columns}\label{sec:tables-advanced-4}

The column specification may fix the padding around a column with an empty at-sign
argument, which removes the space at one end of the table. Repeated columns are written
with a star argument, which expands to a run of identical columns.

\begin{center}
\begin{tabular}{@{}*{3}{c}@{}}
Left & Middle & Right \\
1 & 2 & 3 \\
4 & 5 & 6 \\
\end{tabular}
\end{center}

Each column keeps its own alignment in every row, so the numbers in the table below
line up on their right edges whatever the width of their labels.

\begin{center}
\begin{tabular}{l c r}
Run & Pages & Seconds \\
Alpha & 12 & 0.8 \\
Beta  & 29 & 1.9 \\
Gamma & 61 & 4.2 \\
\end{tabular}
\end{center}

Extra vertical space between rows is written after a row break with an optional length,
as in \texttt{[2pt]}. The table below uses that spacing to separate groups of rows.

\begin{center}
\begin{tabular}{l r}
Operation & Cost \\[2pt]
Read      & 3 \\
Write     & 7 \\[2pt]
Reference & 11 \\
\end{tabular}
\end{center}

Tables with many rows are common in reports, and the author who sets them by hand often
spends more time on alignment than on content. The tabular environment reduces that
work to a column specification and a list of rows. Each row ends with a double
backslash, and each cell is separated from its neighbour by an ampersand. The
alignment of each column is taken from the specification, so a change of one letter
moves the whole column without any other edit.

The tables so far show the alignment letters, the padding that the at-sign notation
controls and the repetition of a column pattern. Horizontal rules, partial rules and
cells that span several columns are set by the same machinery.

\subsection{Rules and spanning cells}\label{sec:tables-rules-4}
Table rules come in three kinds: a full rule across every column, a partial rule
across some of them, and a doubled rule made of two full rules in a row. A heading can
span all the columns of the table, and a single cell can span two of them.

\begin{center}
\begin{tabular}{|l|c|r|}
\hline
\multicolumn{3}{|c|}{\bfseries Measured costs} \\
\hline\hline
Phase & Calls & Time \\
\cline{2-3}
Scan    & 120 & 3.5 \\
Expand  & 940 & 12.0 \\
\multicolumn{2}{|l|}{Total, both phases} & 15.5 \\
\hline
\end{tabular}
\end{center}

\section{Floats}\label{sec:floats-4}

Floating material such as tables and figures is moved by the page builder to a position
where it fits. The placement specifier suggests where the float may go: \texttt{h} for
here, \texttt{t} for the top of a page, \texttt{b} for the bottom and \texttt{p} for a
page of floats. The author can give several letters at once, as in \texttt{[htbp]}, and
the page builder tries them in turn.

Table~\ref{tab:results-4} appears before the text that discusses it, so the reader meets
the reference first and the float later. The float is placed here by its own
specifier.

\begin{table}[h]
\centering
\begin{tabular}{l c r}
Input size & Lines & Pages \\
Short      & 40    & 1     \\
Medium     & 400   & 12    \\
Long       & 4000  & 120   \\
\end{tabular}
\caption{Pages produced for three input sizes.}\label{tab:results-4}
\end{table}

The results in Table~\ref{tab:results-4} grow with the input, and the page counts follow
the line counts closely. The next float is a figure with a caption long enough to wrap
onto a second line, and it is placed at the top of the page by its specifier.

\begin{figure}[t]
\centering
\mbox{\textsf{Input}\quad$\rightarrow$\quad\texttt{Parser}\quad$\rightarrow$\quad\textsf{Layout}\quad$\rightarrow$\quad\textsf{PDF}}
\caption{The stages of the native pipeline, from the characters of the source file to the
bytes of the output file, shown as a chain of boxes in the order in which the data flows
through them.}\label{fig:pipeline-4}
\end{figure}

Figure~\ref{fig:pipeline-4} shows the chain of stages that the native path runs in a
single pass. The figure caption is set as a paragraph because it runs over two lines.
A table placed at the bottom of a page is specified with \texttt{[b]}, as the next
example shows.

\begin{table}[b]
\centering
\begin{tabular}{c c}
Feature & State \\
Lists        & done \\
Tables       & done \\
\end{tabular}
\caption{Feature status at the time of writing.}
\end{table}

A float with both the top and the bottom permitted is written \texttt{[tb]}, and a float
that may go anywhere in the sequence is written \texttt{[htbp]}. The next figure uses
the latter specifier and a label that the text refers to from a later section.

\begin{figure}[htbp]
\centering
\mbox{\texttt{figure}} \quad \mbox{Text placed in a box}
\caption{A short caption.}\label{fig:box-4}
\end{figure}

Figure~\ref{fig:box-4} is referred to here, after the float itself. The references to the
tables and figures of this section are resolved in the second pass, which is the reason
that the reader sees the correct numbers even for references that come before their
floats.

A float is a block of material that the page builder may move away from the place where
it appears in the source. Its purpose is to keep large items, such as a table of
results or a diagram, from splitting across a page break or from leaving a large gap
at the foot of a page. The page builder reads the placement specifier of the float and
considers each candidate position in turn. If no candidate fits, the float waits for
the next page, and the material that follows it continues to fill the current page.

A caption belongs to its float and is numbered with it. The number of a table and the
number of a figure are kept in separate counters, so that a document may contain a
table 1 and a figure 1 without conflict. The label placed after the caption records the
number, and the reference in the text picks up the number from the label when the
document is compiled. This is the reason that a reference may come before its float in
the source, as the references to the tables above do.

\section{Mathematics}\label{sec:math-4}

Mathematical notation is set in math mode, which uses its own fonts and spacing rules.
Inline formulae are written between dollar signs. The quantity $\frac{a+b}{c-d}$ is a
fraction, and $\sqrt{x^2+1}$ is a square root. A superscript such as $x^{2}$ and a
subscript such as $a_{i}$ are common, and both can appear together, as in $x_{i}^{2}$.

Greek letters are typed by name, so $\alpha$, $\beta$ and $\gamma$ are available, and so
is the capital $\Gamma$. The sum $\sum_{k=1}^{n} k$ and the integral $\int_0^1 x\,dx$
appear inline with their limits set beside them, as is customary in running text. The
limit $\lim_{n\to\infty} \frac{1}{n} = 0$ also appears inline.

Sets are typeset with blackboard and calligraphic letters. The set of real numbers is
$\mathbb{R}$, and a function that bounds another is written $f(n) \in \mathcal{O}(n^2)$.
A vector is written in bold, as $\mathbf{v}$, and its length is $\|\mathbf{v}\|$. The
parentheses in $\left( \frac{x}{y} \right)$ grow with the fraction they enclose.

Displayed formulae are set on their own line, centred, with a larger style than inline
formulae. The formula below gives the sum of the first $n$ integers.
\[
  \sum_{k=1}^{n} k = \frac{n(n+1)}{2}
\]

Numbered equations carry a label, so that the text can refer to them by number. The
first equation of this section is the definition of the Gaussian integral, and the
second is the quadratic formula, which is used in Section~\ref{sec:lists-4}.

\begin{equation}\label{eq:gauss-4}
  \int_{-\infty}^{\infty} e^{-x^{2}}\,dx = \sqrt{\pi}
\end{equation}

\begin{equation}\label{eq:quadratic-4}
  x = \frac{-b \pm \sqrt{b^{2}-4ac}}{2a}
\end{equation}

Equation~\ref{eq:gauss-4} gives the Gaussian integral, and
Equation~\ref{eq:quadratic-4} gives the roots of a quadratic polynomial. The coefficients
$a$, $b$ and $c$ in Equation~\ref{eq:quadratic-4} are real numbers with $a \neq 0$.

A last display shows a sum over a product, with a subscript that is itself a formula:
\[
  \prod_{i=1}^{m} \left( 1 + \frac{1}{i^{2}} \right) \leq e^{\pi^{2}/6}
\]

Mathematical formulae are read as a sequence of symbols with a structure, and the layout
rules give each symbol a size and a position. Inline formulae are set in the text style,
which keeps the lines of a paragraph at a uniform height, while displayed formulae are
set in the display style, which allows larger operators and fractions with larger
numerators and denominators. The reader sees the difference in the sums and products
above, where the limits sit above and below the symbol in the display, and beside it in
the running text.

The quadratic formula in Equation~\ref{eq:quadratic-4} is a familiar example of a
fraction with a root in its numerator. The square root of the discriminant
$b^{2}-4ac$ is set with a radical sign that grows with its contents, and the
fraction line extends beneath the whole numerator. The parentheses of the next example
are sized in the same way, so that the reader can see the structure at a glance.

\subsection{Displayed derivations}\label{sec:derivations-4}

A derivation shows a chain of equalities, each following from the one before. Each line
of the chain is written as a separate display, and the text between the displays gives
the reason for each step. The first derivation below expands the square of a binomial,
and the second uses the result to evaluate a finite sum of squares.
\[
  (a+b)^{2} = a^{2} + 2ab + b^{2}
\]

Applying the same expansion to the terms of a sum gives a closed form for the sum of the
squares of the first $n$ integers, which is the second formula of this subsection:
\[
  \sum_{k=1}^{n} k^{2} = \frac{n(n+1)(2n+1)}{6}
\]

The formula holds for every positive integer $n$, and it can be proved by induction on
$n$, as the reader may check for the first few values. Each step of the induction
uses the identity above and the fact that the sum for $n+1$ is the sum for $n$ with one
more term added.

The third display is a limit that defines the constant $e$. It is set on its own line
because it is a result that the text refers to, and the reader may wish to compare it
with the bound given in the last display of the previous section:
\[
  e = \lim_{n\to\infty} \left( 1 + \frac{1}{n} \right)^{n}
\]

Derivations that are long are best broken into several displays, each short enough to
fit on one line of the text block, so that the reader can follow the argument without
turning the page.

\section{Bibliography and citations}\label{sec:refs-4}

The native path follows the algorithms of the typesetting system closely, and the
descriptions of those algorithms are found in the standard references. The line
breaking method used by the paragraph builder is described by Knuth and
Plass~\cite{knuth81-4}, while the page builder and the float placement are described in
the reference manual~\cite{knuth86-4}. The layout of mathematical formulae follows the
rules of the appendix on mathematics in that same manual. The chain of stages in
Figure~\ref{fig:pipeline-4} is a common design for compilers in general, as Aho and
colleagues describe~\cite{aho06-4}.

\begin{thebibliography}{9}
\bibitem{knuth81-4}
  Donald E.~Knuth and Michael F.~Plass.
  \emph{Breaking paragraphs into lines}.
  Software: Practice and Experience, 11(11), 1981.

\bibitem{knuth86-4}
  Donald E.~Knuth.
  \emph{The \TeX book}.
  Addison-Wesley, 1986.

\bibitem{aho06-4}
  Alfred V.~Aho, Monica S.~Lam, Ravi Sethi and Jeffrey D.~Ullman.
  \emph{Compilers: Principles, Techniques, and Tools}.
  Addison-Wesley, second edition, 2006.
\end{thebibliography}



\section{Running text}\label{sec:text-5}

Typesetting is the business of turning a stream of characters into a page. A document
author writes paragraphs of plain prose, and the typesetter decides where lines break,
how much space separates words, and where the page ends. The quality of the result
depends on many small decisions, most of which the author never sees. This section
collects the inline devices that an author reaches for while writing prose: emphasis,
font families, font shapes, accents and punctuation.

Emphasis is usually marked with \emph{italic} text, which in an upright paragraph
reverses to roman when nested inside an italic context. Strong statements can be set in
\textbf{bold}, while \textit{explicit italic} and \textsc{small capitals} are useful for
names and terms that deserve a visual distinction without strong emphasis. Source code
and file names belong in \texttt{typewriter type}, and interface labels read best in
\textsf{sans serif}.

Font declarations act on the rest of the current group. The sentence {\bfseries this
phrase is set in bold by a declaration}, and {\itshape this one by an italic
declaration}, so that the surrounding text is unaffected. The older forms \bf and \it
work the same way inside braces, as in {\bf old-style bold} and {\it old-style italic}.

Accents are part of ordinary text. The word caf\'e takes an acute accent, the name
G\"odel takes an umlaut, and the phrase na\"\i ve reads as it should. Quotations are
typed as ``double quotes'' and the apostrophe in don't is a plain character, while
the dashes -- en and --- em -- have their own meanings.

Sentences that trail away end with \ldots{} and the typesetter spaces them as a unit.
The name of the system is written as \TeX{} in the running text and \LaTeX{} in the
document class that builds on it. Non-breaking spaces such as Figure~1 and Section~2
keep a label with its number, and a footnote gives a remark that would interrupt the
argument\footnote{A footnote is set at the foot of the page where its reference
appears, in a smaller size, and it is numbered in sequence with the other footnotes of
the document.}.

Explicit line breaks are useful in addresses and short verse:\\
the first line ends at the double backslash,\\
and the second line starts afresh.

Horizontal and vertical space can be inserted by hand. A paragraph that begins with
\noindent is not indented, and an explicit horizontal gap of \hspace{1em} separates
the two words in this sentence. Vertical space is more common between blocks, so the
next paragraph follows after \vspace{6pt} of extra room.

\noindent The paragraph above was started with the command \texttt{noindent}, so it
has no first-line indentation. The paragraphs that follow are indented as usual, which
is the default behaviour of the article class.

Ordinary paragraphs are separated by a blank line. The line breaking algorithm looks at
the whole paragraph before it commits to any break, which is why a change in one word
can move words several lines further down. The author rarely needs to intervene, but
the mechanisms exist, and the rest of this document uses them where they matter.

\subsection{Font declarations in context}\label{sec:fonts-5}

A declaration such as \textsf{bfseries} changes the font until the end of the enclosing
group. The next paragraph shows how groups limit the effect of such declarations, and
how the text returns to its original shape when the group closes.

{\bfseries Inside this group the text is bold, and a word such as \emph{this} still
changes to italic within bold text.} Outside the group the text returns to roman, as
the reader can see in this sentence. The same holds for {\itshape the slanted shape}
and for {\sffamily the sans serif family}, which all end where their braces end.

\subsection{Spacing and hyphenation}\label{sec:spacing-5}

Word spacing in a justified line is not constant. The stretch that a line receives is
shared among its interword glue, and the amount of stretch depends on how much room the
line has after its words are placed. A line with few words and a large gap is
noticeably loose, and a line with many words and little room may be tight. The
typesetter tries to keep all lines of a paragraph as even as possible, and it accepts
a slightly loose line when the alternative is a worse paragraph as a whole.

Long words that do not fit at the end of a line are hyphenated at points chosen from a
table of patterns. The table is built from word lists and the rules that govern where
a word may be divided, and it lets the typesetter break words such as typesetting,
hyphenation, algorithm and reproducibility at reasonable places. A hyphen at the end of
a line is a cost, and two such hyphens in a row are a larger cost, so the algorithm
avoids them where it can.

Some words must never be divided, and the author can say so by writing the word inside
a group that forbids breaks, or by joining the parts of a name with a tie. The name
Ludwig~van Beethoven keeps its two parts together, and a unit such as 12~pt does the
same for the number and the unit.

Sentences that end with a capital letter, such as the initials in J.~R.\ Smith, have a
different space after them, and the author must mark an end of sentence when the
capital is not the end of a sentence. The rule is simple, but it is easy to forget, and
the document below uses a backslash space in one case to show how it is written.

\section{Lists}\label{sec:lists-5}

Lists organise material that has no natural sentence structure. The three list
environments of the article class differ in how they mark their items and how they
number them. The first is the unordered list, which uses a bullet for each item.

\begin{itemize}
  \item The first item of a bulleted list is a plain sentence.
  \item The second item is longer and contains an \emph{emphasised} word, so that the
        wrapped line is indented to the same position as the text of the item.
  \item A third item may itself hold a sublist:
    \begin{itemize}
      \item the nested list uses a dash as its marker,
      \item and its items are indented further than the parent items.
    \end{itemize}
\end{itemize}

Numbered lists are used when the order of the items matters, for example in a
procedure. The following enumeration is nested two levels deep, and its first item
carries a label so that the text can refer to it later. The reference to that item
appears in Section~\ref{sec:tables-5}, which is written after this list.

\begin{enumerate}
  \item\label{item:first-5} Obtain the source files of the document and the packages that
        it loads.
  \item Compile the document once on the native path, which reads the auxiliary
        information of the previous run, if there is one.
        \begin{enumerate}
          \item The first sub-item describes a check that the author can make by hand.
          \item The second sub-item describes what happens when the check fails.
        \end{enumerate}
  \item Compare the result against the reference output and report any page that
        differs.
\end{enumerate}

The first step of the procedure is item~\ref{item:first-5} in the list above, and the
reader who wants the reasons behind the procedure will find them in
Section~\ref{sec:displays-5}.

Definition lists give a term followed by its explanation. The description environment
sets each term in bold and indents the explanation beneath it.

\begin{description}
  \item[Kerning] An adjustment of the space between two particular letters, such as
        the pair A and V, so that the shapes look evenly spaced.
  \item[Ligature] A single glyph that replaces a sequence of letters, such as the
        fi ligature in the word first.
  \item[Leading] The distance from one baseline to the next, measured in points.
\end{description}

The three kinds of list above are the core of most structured writing. An author who
needs a list inside a list simply nests the environment, and each level takes its own
label and its own indentation. The nesting depth in the article class is limited only
by the page, although a list nested more than three levels deep is usually a sign that
the material wants a table or a section instead.

Lists also interact with the page. A list is a sequence of paragraphs with extra space
between items, and the page builder may break the page between two items, as it may
between two paragraphs. The space at the top and bottom of a list is controlled by the
document class, so that a list opens and closes with the same visual distance as a
displayed equation.

\section{Displayed text}\label{sec:displays-5}

Some text is set apart from the body of the document to show that it is a quotation, a
centred title, or a block of literal material. Each environment has its own margins and
its own vertical spacing, and the environments below show the common cases.

\begin{quote}
A quotation environment sets its text with extra margin on both sides and without
paragraph indentation. It is intended for a short passage that the author wishes to
present as a separate block of prose, without the extra emphasis of a block quotation.
\end{quote}

\begin{quotation}
The quotation environment differs from the quote environment in that it indents the
first line of each of its paragraphs. A long passage that runs over several paragraphs
therefore keeps its internal structure, as a reader expects from a transcript of a
speech or a letter. This quotation has two paragraphs, so the second one can be seen to
start with an indent.

The second paragraph of the quotation shows the indentation of a continuing paragraph.
\end{quotation}

\begin{center}
Centred text on a line of its own\\
keeps each of its lines in the middle of the text block.
\end{center}

\begin{flushleft}
Flush left text keeps the left margin straight and lets the right margin
ragged, which suits material such as addresses and minutes of a meeting.
\end{flushleft}

\begin{flushright}
Flush right text keeps the right margin straight, which is seldom useful for
prose but is useful for a signature line.
\end{flushright}

Verbatim text is typeset exactly as it is typed, with spaces, line breaks and special
characters such as the percent sign and the braces kept as they stand:

\begin{verbatim}
  for i in 1..n {
      total += a[i] * b[i];   % dot product, not a comment
  }
\end{verbatim}

The following block uses a centring declaration inside a group, which keeps the
alignment local to the braces, and then a ragged right declaration on the next block.

{\centering A centred block made with a declaration, which ends at the closing
brace and leaves the surrounding text alone.\par}

{\raggedright A block whose lines are not stretched to the full measure, so the
right margin is ragged. The declaration ends with the group, and the next paragraph
returns to the usual justified setting of the text.\par}

Paragraphs of ordinary text after these blocks are justified again, as the reader can
see here.

\section{Tables}\label{sec:tables-5}

Tables set data in rows and columns. The tabular environment takes a column
specification, and each column is either left, centred or right aligned. Vertical bars
draw rules between columns. Section~\ref{sec:displays-5} explains the centring that surrounds each table in
this section.

\subsection{Basic alignment and rules}\label{sec:tables-basic-5}

\begin{center}
\begin{tabular}{l c r}
Method & Passes & Time (ms) \\
Greedy fill      & 1 & 0.4 \\
Knuth--Plass     & 2 & 3.1 \\
Optimal fit      & 3 & 4.7 \\
\end{tabular}
\end{center}

\begin{center}
\begin{tabular}{|l|c|r|}
Language & Typing & Speed \\
C        & static  & fast   \\
Python   & dynamic & slow   \\
Rust     & static  & fast   \\
\end{tabular}
\end{center}

\subsection{Padding and repeated columns}\label{sec:tables-advanced-5}

The column specification may fix the padding around a column with an empty at-sign
argument, which removes the space at one end of the table. Repeated columns are written
with a star argument, which expands to a run of identical columns.

\begin{center}
\begin{tabular}{@{}*{3}{c}@{}}
Left & Middle & Right \\
1 & 2 & 3 \\
4 & 5 & 6 \\
\end{tabular}
\end{center}

Each column keeps its own alignment in every row, so the numbers in the table below
line up on their right edges whatever the width of their labels.

\begin{center}
\begin{tabular}{l c r}
Run & Pages & Seconds \\
Alpha & 12 & 0.8 \\
Beta  & 29 & 1.9 \\
Gamma & 61 & 4.2 \\
\end{tabular}
\end{center}

Extra vertical space between rows is written after a row break with an optional length,
as in \texttt{[2pt]}. The table below uses that spacing to separate groups of rows.

\begin{center}
\begin{tabular}{l r}
Operation & Cost \\[2pt]
Read      & 3 \\
Write     & 7 \\[2pt]
Reference & 11 \\
\end{tabular}
\end{center}

Tables with many rows are common in reports, and the author who sets them by hand often
spends more time on alignment than on content. The tabular environment reduces that
work to a column specification and a list of rows. Each row ends with a double
backslash, and each cell is separated from its neighbour by an ampersand. The
alignment of each column is taken from the specification, so a change of one letter
moves the whole column without any other edit.

The tables so far show the alignment letters, the padding that the at-sign notation
controls and the repetition of a column pattern. Horizontal rules, partial rules and
cells that span several columns are set by the same machinery.

\subsection{Rules and spanning cells}\label{sec:tables-rules-5}
Table rules come in three kinds: a full rule across every column, a partial rule
across some of them, and a doubled rule made of two full rules in a row. A heading can
span all the columns of the table, and a single cell can span two of them.

\begin{center}
\begin{tabular}{|l|c|r|}
\hline
\multicolumn{3}{|c|}{\bfseries Measured costs} \\
\hline\hline
Phase & Calls & Time \\
\cline{2-3}
Scan    & 120 & 3.5 \\
Expand  & 940 & 12.0 \\
\multicolumn{2}{|l|}{Total, both phases} & 15.5 \\
\hline
\end{tabular}
\end{center}

\section{Floats}\label{sec:floats-5}

Floating material such as tables and figures is moved by the page builder to a position
where it fits. The placement specifier suggests where the float may go: \texttt{h} for
here, \texttt{t} for the top of a page, \texttt{b} for the bottom and \texttt{p} for a
page of floats. The author can give several letters at once, as in \texttt{[htbp]}, and
the page builder tries them in turn.

Table~\ref{tab:results-5} appears before the text that discusses it, so the reader meets
the reference first and the float later. The float is placed here by its own
specifier.

\begin{table}[h]
\centering
\begin{tabular}{l c r}
Input size & Lines & Pages \\
Short      & 40    & 1     \\
Medium     & 400   & 12    \\
Long       & 4000  & 120   \\
\end{tabular}
\caption{Pages produced for three input sizes.}\label{tab:results-5}
\end{table}

The results in Table~\ref{tab:results-5} grow with the input, and the page counts follow
the line counts closely. The next float is a figure with a caption long enough to wrap
onto a second line, and it is placed at the top of the page by its specifier.

\begin{figure}[t]
\centering
\mbox{\textsf{Input}\quad$\rightarrow$\quad\texttt{Parser}\quad$\rightarrow$\quad\textsf{Layout}\quad$\rightarrow$\quad\textsf{PDF}}
\caption{The stages of the native pipeline, from the characters of the source file to the
bytes of the output file, shown as a chain of boxes in the order in which the data flows
through them.}\label{fig:pipeline-5}
\end{figure}

Figure~\ref{fig:pipeline-5} shows the chain of stages that the native path runs in a
single pass. The figure caption is set as a paragraph because it runs over two lines.
A table placed at the bottom of a page is specified with \texttt{[b]}, as the next
example shows.

\begin{table}[b]
\centering
\begin{tabular}{c c}
Feature & State \\
Lists        & done \\
Tables       & done \\
\end{tabular}
\caption{Feature status at the time of writing.}
\end{table}

A float with both the top and the bottom permitted is written \texttt{[tb]}, and a float
that may go anywhere in the sequence is written \texttt{[htbp]}. The next figure uses
the latter specifier and a label that the text refers to from a later section.

\begin{figure}[htbp]
\centering
\mbox{\texttt{figure}} \quad \mbox{Text placed in a box}
\caption{A short caption.}\label{fig:box-5}
\end{figure}

Figure~\ref{fig:box-5} is referred to here, after the float itself. The references to the
tables and figures of this section are resolved in the second pass, which is the reason
that the reader sees the correct numbers even for references that come before their
floats.

A float is a block of material that the page builder may move away from the place where
it appears in the source. Its purpose is to keep large items, such as a table of
results or a diagram, from splitting across a page break or from leaving a large gap
at the foot of a page. The page builder reads the placement specifier of the float and
considers each candidate position in turn. If no candidate fits, the float waits for
the next page, and the material that follows it continues to fill the current page.

A caption belongs to its float and is numbered with it. The number of a table and the
number of a figure are kept in separate counters, so that a document may contain a
table 1 and a figure 1 without conflict. The label placed after the caption records the
number, and the reference in the text picks up the number from the label when the
document is compiled. This is the reason that a reference may come before its float in
the source, as the references to the tables above do.

\section{Mathematics}\label{sec:math-5}

Mathematical notation is set in math mode, which uses its own fonts and spacing rules.
Inline formulae are written between dollar signs. The quantity $\frac{a+b}{c-d}$ is a
fraction, and $\sqrt{x^2+1}$ is a square root. A superscript such as $x^{2}$ and a
subscript such as $a_{i}$ are common, and both can appear together, as in $x_{i}^{2}$.

Greek letters are typed by name, so $\alpha$, $\beta$ and $\gamma$ are available, and so
is the capital $\Gamma$. The sum $\sum_{k=1}^{n} k$ and the integral $\int_0^1 x\,dx$
appear inline with their limits set beside them, as is customary in running text. The
limit $\lim_{n\to\infty} \frac{1}{n} = 0$ also appears inline.

Sets are typeset with blackboard and calligraphic letters. The set of real numbers is
$\mathbb{R}$, and a function that bounds another is written $f(n) \in \mathcal{O}(n^2)$.
A vector is written in bold, as $\mathbf{v}$, and its length is $\|\mathbf{v}\|$. The
parentheses in $\left( \frac{x}{y} \right)$ grow with the fraction they enclose.

Displayed formulae are set on their own line, centred, with a larger style than inline
formulae. The formula below gives the sum of the first $n$ integers.
\[
  \sum_{k=1}^{n} k = \frac{n(n+1)}{2}
\]

Numbered equations carry a label, so that the text can refer to them by number. The
first equation of this section is the definition of the Gaussian integral, and the
second is the quadratic formula, which is used in Section~\ref{sec:lists-5}.

\begin{equation}\label{eq:gauss-5}
  \int_{-\infty}^{\infty} e^{-x^{2}}\,dx = \sqrt{\pi}
\end{equation}

\begin{equation}\label{eq:quadratic-5}
  x = \frac{-b \pm \sqrt{b^{2}-4ac}}{2a}
\end{equation}

Equation~\ref{eq:gauss-5} gives the Gaussian integral, and
Equation~\ref{eq:quadratic-5} gives the roots of a quadratic polynomial. The coefficients
$a$, $b$ and $c$ in Equation~\ref{eq:quadratic-5} are real numbers with $a \neq 0$.

A last display shows a sum over a product, with a subscript that is itself a formula:
\[
  \prod_{i=1}^{m} \left( 1 + \frac{1}{i^{2}} \right) \leq e^{\pi^{2}/6}
\]

Mathematical formulae are read as a sequence of symbols with a structure, and the layout
rules give each symbol a size and a position. Inline formulae are set in the text style,
which keeps the lines of a paragraph at a uniform height, while displayed formulae are
set in the display style, which allows larger operators and fractions with larger
numerators and denominators. The reader sees the difference in the sums and products
above, where the limits sit above and below the symbol in the display, and beside it in
the running text.

The quadratic formula in Equation~\ref{eq:quadratic-5} is a familiar example of a
fraction with a root in its numerator. The square root of the discriminant
$b^{2}-4ac$ is set with a radical sign that grows with its contents, and the
fraction line extends beneath the whole numerator. The parentheses of the next example
are sized in the same way, so that the reader can see the structure at a glance.

\subsection{Displayed derivations}\label{sec:derivations-5}

A derivation shows a chain of equalities, each following from the one before. Each line
of the chain is written as a separate display, and the text between the displays gives
the reason for each step. The first derivation below expands the square of a binomial,
and the second uses the result to evaluate a finite sum of squares.
\[
  (a+b)^{2} = a^{2} + 2ab + b^{2}
\]

Applying the same expansion to the terms of a sum gives a closed form for the sum of the
squares of the first $n$ integers, which is the second formula of this subsection:
\[
  \sum_{k=1}^{n} k^{2} = \frac{n(n+1)(2n+1)}{6}
\]

The formula holds for every positive integer $n$, and it can be proved by induction on
$n$, as the reader may check for the first few values. Each step of the induction
uses the identity above and the fact that the sum for $n+1$ is the sum for $n$ with one
more term added.

The third display is a limit that defines the constant $e$. It is set on its own line
because it is a result that the text refers to, and the reader may wish to compare it
with the bound given in the last display of the previous section:
\[
  e = \lim_{n\to\infty} \left( 1 + \frac{1}{n} \right)^{n}
\]

Derivations that are long are best broken into several displays, each short enough to
fit on one line of the text block, so that the reader can follow the argument without
turning the page.

\section{Bibliography and citations}\label{sec:refs-5}

The native path follows the algorithms of the typesetting system closely, and the
descriptions of those algorithms are found in the standard references. The line
breaking method used by the paragraph builder is described by Knuth and
Plass~\cite{knuth81-5}, while the page builder and the float placement are described in
the reference manual~\cite{knuth86-5}. The layout of mathematical formulae follows the
rules of the appendix on mathematics in that same manual. The chain of stages in
Figure~\ref{fig:pipeline-5} is a common design for compilers in general, as Aho and
colleagues describe~\cite{aho06-5}.

\begin{thebibliography}{9}
\bibitem{knuth81-5}
  Donald E.~Knuth and Michael F.~Plass.
  \emph{Breaking paragraphs into lines}.
  Software: Practice and Experience, 11(11), 1981.

\bibitem{knuth86-5}
  Donald E.~Knuth.
  \emph{The \TeX book}.
  Addison-Wesley, 1986.

\bibitem{aho06-5}
  Alfred V.~Aho, Monica S.~Lam, Ravi Sethi and Jeffrey D.~Ullman.
  \emph{Compilers: Principles, Techniques, and Tools}.
  Addison-Wesley, second edition, 2006.
\end{thebibliography}



\section{Running text}\label{sec:text-6}

Typesetting is the business of turning a stream of characters into a page. A document
author writes paragraphs of plain prose, and the typesetter decides where lines break,
how much space separates words, and where the page ends. The quality of the result
depends on many small decisions, most of which the author never sees. This section
collects the inline devices that an author reaches for while writing prose: emphasis,
font families, font shapes, accents and punctuation.

Emphasis is usually marked with \emph{italic} text, which in an upright paragraph
reverses to roman when nested inside an italic context. Strong statements can be set in
\textbf{bold}, while \textit{explicit italic} and \textsc{small capitals} are useful for
names and terms that deserve a visual distinction without strong emphasis. Source code
and file names belong in \texttt{typewriter type}, and interface labels read best in
\textsf{sans serif}.

Font declarations act on the rest of the current group. The sentence {\bfseries this
phrase is set in bold by a declaration}, and {\itshape this one by an italic
declaration}, so that the surrounding text is unaffected. The older forms \bf and \it
work the same way inside braces, as in {\bf old-style bold} and {\it old-style italic}.

Accents are part of ordinary text. The word caf\'e takes an acute accent, the name
G\"odel takes an umlaut, and the phrase na\"\i ve reads as it should. Quotations are
typed as ``double quotes'' and the apostrophe in don't is a plain character, while
the dashes -- en and --- em -- have their own meanings.

Sentences that trail away end with \ldots{} and the typesetter spaces them as a unit.
The name of the system is written as \TeX{} in the running text and \LaTeX{} in the
document class that builds on it. Non-breaking spaces such as Figure~1 and Section~2
keep a label with its number, and a footnote gives a remark that would interrupt the
argument\footnote{A footnote is set at the foot of the page where its reference
appears, in a smaller size, and it is numbered in sequence with the other footnotes of
the document.}.

Explicit line breaks are useful in addresses and short verse:\\
the first line ends at the double backslash,\\
and the second line starts afresh.

Horizontal and vertical space can be inserted by hand. A paragraph that begins with
\noindent is not indented, and an explicit horizontal gap of \hspace{1em} separates
the two words in this sentence. Vertical space is more common between blocks, so the
next paragraph follows after \vspace{6pt} of extra room.

\noindent The paragraph above was started with the command \texttt{noindent}, so it
has no first-line indentation. The paragraphs that follow are indented as usual, which
is the default behaviour of the article class.

Ordinary paragraphs are separated by a blank line. The line breaking algorithm looks at
the whole paragraph before it commits to any break, which is why a change in one word
can move words several lines further down. The author rarely needs to intervene, but
the mechanisms exist, and the rest of this document uses them where they matter.

\subsection{Font declarations in context}\label{sec:fonts-6}

A declaration such as \textsf{bfseries} changes the font until the end of the enclosing
group. The next paragraph shows how groups limit the effect of such declarations, and
how the text returns to its original shape when the group closes.

{\bfseries Inside this group the text is bold, and a word such as \emph{this} still
changes to italic within bold text.} Outside the group the text returns to roman, as
the reader can see in this sentence. The same holds for {\itshape the slanted shape}
and for {\sffamily the sans serif family}, which all end where their braces end.

\subsection{Spacing and hyphenation}\label{sec:spacing-6}

Word spacing in a justified line is not constant. The stretch that a line receives is
shared among its interword glue, and the amount of stretch depends on how much room the
line has after its words are placed. A line with few words and a large gap is
noticeably loose, and a line with many words and little room may be tight. The
typesetter tries to keep all lines of a paragraph as even as possible, and it accepts
a slightly loose line when the alternative is a worse paragraph as a whole.

Long words that do not fit at the end of a line are hyphenated at points chosen from a
table of patterns. The table is built from word lists and the rules that govern where
a word may be divided, and it lets the typesetter break words such as typesetting,
hyphenation, algorithm and reproducibility at reasonable places. A hyphen at the end of
a line is a cost, and two such hyphens in a row are a larger cost, so the algorithm
avoids them where it can.

Some words must never be divided, and the author can say so by writing the word inside
a group that forbids breaks, or by joining the parts of a name with a tie. The name
Ludwig~van Beethoven keeps its two parts together, and a unit such as 12~pt does the
same for the number and the unit.

Sentences that end with a capital letter, such as the initials in J.~R.\ Smith, have a
different space after them, and the author must mark an end of sentence when the
capital is not the end of a sentence. The rule is simple, but it is easy to forget, and
the document below uses a backslash space in one case to show how it is written.

\section{Lists}\label{sec:lists-6}

Lists organise material that has no natural sentence structure. The three list
environments of the article class differ in how they mark their items and how they
number them. The first is the unordered list, which uses a bullet for each item.

\begin{itemize}
  \item The first item of a bulleted list is a plain sentence.
  \item The second item is longer and contains an \emph{emphasised} word, so that the
        wrapped line is indented to the same position as the text of the item.
  \item A third item may itself hold a sublist:
    \begin{itemize}
      \item the nested list uses a dash as its marker,
      \item and its items are indented further than the parent items.
    \end{itemize}
\end{itemize}

Numbered lists are used when the order of the items matters, for example in a
procedure. The following enumeration is nested two levels deep, and its first item
carries a label so that the text can refer to it later. The reference to that item
appears in Section~\ref{sec:tables-6}, which is written after this list.

\begin{enumerate}
  \item\label{item:first-6} Obtain the source files of the document and the packages that
        it loads.
  \item Compile the document once on the native path, which reads the auxiliary
        information of the previous run, if there is one.
        \begin{enumerate}
          \item The first sub-item describes a check that the author can make by hand.
          \item The second sub-item describes what happens when the check fails.
        \end{enumerate}
  \item Compare the result against the reference output and report any page that
        differs.
\end{enumerate}

The first step of the procedure is item~\ref{item:first-6} in the list above, and the
reader who wants the reasons behind the procedure will find them in
Section~\ref{sec:displays-6}.

Definition lists give a term followed by its explanation. The description environment
sets each term in bold and indents the explanation beneath it.

\begin{description}
  \item[Kerning] An adjustment of the space between two particular letters, such as
        the pair A and V, so that the shapes look evenly spaced.
  \item[Ligature] A single glyph that replaces a sequence of letters, such as the
        fi ligature in the word first.
  \item[Leading] The distance from one baseline to the next, measured in points.
\end{description}

The three kinds of list above are the core of most structured writing. An author who
needs a list inside a list simply nests the environment, and each level takes its own
label and its own indentation. The nesting depth in the article class is limited only
by the page, although a list nested more than three levels deep is usually a sign that
the material wants a table or a section instead.

Lists also interact with the page. A list is a sequence of paragraphs with extra space
between items, and the page builder may break the page between two items, as it may
between two paragraphs. The space at the top and bottom of a list is controlled by the
document class, so that a list opens and closes with the same visual distance as a
displayed equation.

\section{Displayed text}\label{sec:displays-6}

Some text is set apart from the body of the document to show that it is a quotation, a
centred title, or a block of literal material. Each environment has its own margins and
its own vertical spacing, and the environments below show the common cases.

\begin{quote}
A quotation environment sets its text with extra margin on both sides and without
paragraph indentation. It is intended for a short passage that the author wishes to
present as a separate block of prose, without the extra emphasis of a block quotation.
\end{quote}

\begin{quotation}
The quotation environment differs from the quote environment in that it indents the
first line of each of its paragraphs. A long passage that runs over several paragraphs
therefore keeps its internal structure, as a reader expects from a transcript of a
speech or a letter. This quotation has two paragraphs, so the second one can be seen to
start with an indent.

The second paragraph of the quotation shows the indentation of a continuing paragraph.
\end{quotation}

\begin{center}
Centred text on a line of its own\\
keeps each of its lines in the middle of the text block.
\end{center}

\begin{flushleft}
Flush left text keeps the left margin straight and lets the right margin
ragged, which suits material such as addresses and minutes of a meeting.
\end{flushleft}

\begin{flushright}
Flush right text keeps the right margin straight, which is seldom useful for
prose but is useful for a signature line.
\end{flushright}

Verbatim text is typeset exactly as it is typed, with spaces, line breaks and special
characters such as the percent sign and the braces kept as they stand:

\begin{verbatim}
  for i in 1..n {
      total += a[i] * b[i];   % dot product, not a comment
  }
\end{verbatim}

The following block uses a centring declaration inside a group, which keeps the
alignment local to the braces, and then a ragged right declaration on the next block.

{\centering A centred block made with a declaration, which ends at the closing
brace and leaves the surrounding text alone.\par}

{\raggedright A block whose lines are not stretched to the full measure, so the
right margin is ragged. The declaration ends with the group, and the next paragraph
returns to the usual justified setting of the text.\par}

Paragraphs of ordinary text after these blocks are justified again, as the reader can
see here.

\section{Tables}\label{sec:tables-6}

Tables set data in rows and columns. The tabular environment takes a column
specification, and each column is either left, centred or right aligned. Vertical bars
draw rules between columns. Section~\ref{sec:displays-6} explains the centring that surrounds each table in
this section.

\subsection{Basic alignment and rules}\label{sec:tables-basic-6}

\begin{center}
\begin{tabular}{l c r}
Method & Passes & Time (ms) \\
Greedy fill      & 1 & 0.4 \\
Knuth--Plass     & 2 & 3.1 \\
Optimal fit      & 3 & 4.7 \\
\end{tabular}
\end{center}

\begin{center}
\begin{tabular}{|l|c|r|}
Language & Typing & Speed \\
C        & static  & fast   \\
Python   & dynamic & slow   \\
Rust     & static  & fast   \\
\end{tabular}
\end{center}

\subsection{Padding and repeated columns}\label{sec:tables-advanced-6}

The column specification may fix the padding around a column with an empty at-sign
argument, which removes the space at one end of the table. Repeated columns are written
with a star argument, which expands to a run of identical columns.

\begin{center}
\begin{tabular}{@{}*{3}{c}@{}}
Left & Middle & Right \\
1 & 2 & 3 \\
4 & 5 & 6 \\
\end{tabular}
\end{center}

Each column keeps its own alignment in every row, so the numbers in the table below
line up on their right edges whatever the width of their labels.

\begin{center}
\begin{tabular}{l c r}
Run & Pages & Seconds \\
Alpha & 12 & 0.8 \\
Beta  & 29 & 1.9 \\
Gamma & 61 & 4.2 \\
\end{tabular}
\end{center}

Extra vertical space between rows is written after a row break with an optional length,
as in \texttt{[2pt]}. The table below uses that spacing to separate groups of rows.

\begin{center}
\begin{tabular}{l r}
Operation & Cost \\[2pt]
Read      & 3 \\
Write     & 7 \\[2pt]
Reference & 11 \\
\end{tabular}
\end{center}

Tables with many rows are common in reports, and the author who sets them by hand often
spends more time on alignment than on content. The tabular environment reduces that
work to a column specification and a list of rows. Each row ends with a double
backslash, and each cell is separated from its neighbour by an ampersand. The
alignment of each column is taken from the specification, so a change of one letter
moves the whole column without any other edit.

The tables so far show the alignment letters, the padding that the at-sign notation
controls and the repetition of a column pattern. Horizontal rules, partial rules and
cells that span several columns are set by the same machinery.

\subsection{Rules and spanning cells}\label{sec:tables-rules-6}
Table rules come in three kinds: a full rule across every column, a partial rule
across some of them, and a doubled rule made of two full rules in a row. A heading can
span all the columns of the table, and a single cell can span two of them.

\begin{center}
\begin{tabular}{|l|c|r|}
\hline
\multicolumn{3}{|c|}{\bfseries Measured costs} \\
\hline\hline
Phase & Calls & Time \\
\cline{2-3}
Scan    & 120 & 3.5 \\
Expand  & 940 & 12.0 \\
\multicolumn{2}{|l|}{Total, both phases} & 15.5 \\
\hline
\end{tabular}
\end{center}

\section{Floats}\label{sec:floats-6}

Floating material such as tables and figures is moved by the page builder to a position
where it fits. The placement specifier suggests where the float may go: \texttt{h} for
here, \texttt{t} for the top of a page, \texttt{b} for the bottom and \texttt{p} for a
page of floats. The author can give several letters at once, as in \texttt{[htbp]}, and
the page builder tries them in turn.

Table~\ref{tab:results-6} appears before the text that discusses it, so the reader meets
the reference first and the float later. The float is placed here by its own
specifier.

\begin{table}[h]
\centering
\begin{tabular}{l c r}
Input size & Lines & Pages \\
Short      & 40    & 1     \\
Medium     & 400   & 12    \\
Long       & 4000  & 120   \\
\end{tabular}
\caption{Pages produced for three input sizes.}\label{tab:results-6}
\end{table}

The results in Table~\ref{tab:results-6} grow with the input, and the page counts follow
the line counts closely. The next float is a figure with a caption long enough to wrap
onto a second line, and it is placed at the top of the page by its specifier.

\begin{figure}[t]
\centering
\mbox{\textsf{Input}\quad$\rightarrow$\quad\texttt{Parser}\quad$\rightarrow$\quad\textsf{Layout}\quad$\rightarrow$\quad\textsf{PDF}}
\caption{The stages of the native pipeline, from the characters of the source file to the
bytes of the output file, shown as a chain of boxes in the order in which the data flows
through them.}\label{fig:pipeline-6}
\end{figure}

Figure~\ref{fig:pipeline-6} shows the chain of stages that the native path runs in a
single pass. The figure caption is set as a paragraph because it runs over two lines.
A table placed at the bottom of a page is specified with \texttt{[b]}, as the next
example shows.

\begin{table}[b]
\centering
\begin{tabular}{c c}
Feature & State \\
Lists        & done \\
Tables       & done \\
\end{tabular}
\caption{Feature status at the time of writing.}
\end{table}

A float with both the top and the bottom permitted is written \texttt{[tb]}, and a float
that may go anywhere in the sequence is written \texttt{[htbp]}. The next figure uses
the latter specifier and a label that the text refers to from a later section.

\begin{figure}[htbp]
\centering
\mbox{\texttt{figure}} \quad \mbox{Text placed in a box}
\caption{A short caption.}\label{fig:box-6}
\end{figure}

Figure~\ref{fig:box-6} is referred to here, after the float itself. The references to the
tables and figures of this section are resolved in the second pass, which is the reason
that the reader sees the correct numbers even for references that come before their
floats.

A float is a block of material that the page builder may move away from the place where
it appears in the source. Its purpose is to keep large items, such as a table of
results or a diagram, from splitting across a page break or from leaving a large gap
at the foot of a page. The page builder reads the placement specifier of the float and
considers each candidate position in turn. If no candidate fits, the float waits for
the next page, and the material that follows it continues to fill the current page.

A caption belongs to its float and is numbered with it. The number of a table and the
number of a figure are kept in separate counters, so that a document may contain a
table 1 and a figure 1 without conflict. The label placed after the caption records the
number, and the reference in the text picks up the number from the label when the
document is compiled. This is the reason that a reference may come before its float in
the source, as the references to the tables above do.

\section{Mathematics}\label{sec:math-6}

Mathematical notation is set in math mode, which uses its own fonts and spacing rules.
Inline formulae are written between dollar signs. The quantity $\frac{a+b}{c-d}$ is a
fraction, and $\sqrt{x^2+1}$ is a square root. A superscript such as $x^{2}$ and a
subscript such as $a_{i}$ are common, and both can appear together, as in $x_{i}^{2}$.

Greek letters are typed by name, so $\alpha$, $\beta$ and $\gamma$ are available, and so
is the capital $\Gamma$. The sum $\sum_{k=1}^{n} k$ and the integral $\int_0^1 x\,dx$
appear inline with their limits set beside them, as is customary in running text. The
limit $\lim_{n\to\infty} \frac{1}{n} = 0$ also appears inline.

Sets are typeset with blackboard and calligraphic letters. The set of real numbers is
$\mathbb{R}$, and a function that bounds another is written $f(n) \in \mathcal{O}(n^2)$.
A vector is written in bold, as $\mathbf{v}$, and its length is $\|\mathbf{v}\|$. The
parentheses in $\left( \frac{x}{y} \right)$ grow with the fraction they enclose.

Displayed formulae are set on their own line, centred, with a larger style than inline
formulae. The formula below gives the sum of the first $n$ integers.
\[
  \sum_{k=1}^{n} k = \frac{n(n+1)}{2}
\]

Numbered equations carry a label, so that the text can refer to them by number. The
first equation of this section is the definition of the Gaussian integral, and the
second is the quadratic formula, which is used in Section~\ref{sec:lists-6}.

\begin{equation}\label{eq:gauss-6}
  \int_{-\infty}^{\infty} e^{-x^{2}}\,dx = \sqrt{\pi}
\end{equation}

\begin{equation}\label{eq:quadratic-6}
  x = \frac{-b \pm \sqrt{b^{2}-4ac}}{2a}
\end{equation}

Equation~\ref{eq:gauss-6} gives the Gaussian integral, and
Equation~\ref{eq:quadratic-6} gives the roots of a quadratic polynomial. The coefficients
$a$, $b$ and $c$ in Equation~\ref{eq:quadratic-6} are real numbers with $a \neq 0$.

A last display shows a sum over a product, with a subscript that is itself a formula:
\[
  \prod_{i=1}^{m} \left( 1 + \frac{1}{i^{2}} \right) \leq e^{\pi^{2}/6}
\]

Mathematical formulae are read as a sequence of symbols with a structure, and the layout
rules give each symbol a size and a position. Inline formulae are set in the text style,
which keeps the lines of a paragraph at a uniform height, while displayed formulae are
set in the display style, which allows larger operators and fractions with larger
numerators and denominators. The reader sees the difference in the sums and products
above, where the limits sit above and below the symbol in the display, and beside it in
the running text.

The quadratic formula in Equation~\ref{eq:quadratic-6} is a familiar example of a
fraction with a root in its numerator. The square root of the discriminant
$b^{2}-4ac$ is set with a radical sign that grows with its contents, and the
fraction line extends beneath the whole numerator. The parentheses of the next example
are sized in the same way, so that the reader can see the structure at a glance.

\subsection{Displayed derivations}\label{sec:derivations-6}

A derivation shows a chain of equalities, each following from the one before. Each line
of the chain is written as a separate display, and the text between the displays gives
the reason for each step. The first derivation below expands the square of a binomial,
and the second uses the result to evaluate a finite sum of squares.
\[
  (a+b)^{2} = a^{2} + 2ab + b^{2}
\]

Applying the same expansion to the terms of a sum gives a closed form for the sum of the
squares of the first $n$ integers, which is the second formula of this subsection:
\[
  \sum_{k=1}^{n} k^{2} = \frac{n(n+1)(2n+1)}{6}
\]

The formula holds for every positive integer $n$, and it can be proved by induction on
$n$, as the reader may check for the first few values. Each step of the induction
uses the identity above and the fact that the sum for $n+1$ is the sum for $n$ with one
more term added.

The third display is a limit that defines the constant $e$. It is set on its own line
because it is a result that the text refers to, and the reader may wish to compare it
with the bound given in the last display of the previous section:
\[
  e = \lim_{n\to\infty} \left( 1 + \frac{1}{n} \right)^{n}
\]

Derivations that are long are best broken into several displays, each short enough to
fit on one line of the text block, so that the reader can follow the argument without
turning the page.

\section{Bibliography and citations}\label{sec:refs-6}

The native path follows the algorithms of the typesetting system closely, and the
descriptions of those algorithms are found in the standard references. The line
breaking method used by the paragraph builder is described by Knuth and
Plass~\cite{knuth81-6}, while the page builder and the float placement are described in
the reference manual~\cite{knuth86-6}. The layout of mathematical formulae follows the
rules of the appendix on mathematics in that same manual. The chain of stages in
Figure~\ref{fig:pipeline-6} is a common design for compilers in general, as Aho and
colleagues describe~\cite{aho06-6}.

\begin{thebibliography}{9}
\bibitem{knuth81-6}
  Donald E.~Knuth and Michael F.~Plass.
  \emph{Breaking paragraphs into lines}.
  Software: Practice and Experience, 11(11), 1981.

\bibitem{knuth86-6}
  Donald E.~Knuth.
  \emph{The \TeX book}.
  Addison-Wesley, 1986.

\bibitem{aho06-6}
  Alfred V.~Aho, Monica S.~Lam, Ravi Sethi and Jeffrey D.~Ullman.
  \emph{Compilers: Principles, Techniques, and Tools}.
  Addison-Wesley, second edition, 2006.
\end{thebibliography}



\section{Running text}\label{sec:text-7}

Typesetting is the business of turning a stream of characters into a page. A document
author writes paragraphs of plain prose, and the typesetter decides where lines break,
how much space separates words, and where the page ends. The quality of the result
depends on many small decisions, most of which the author never sees. This section
collects the inline devices that an author reaches for while writing prose: emphasis,
font families, font shapes, accents and punctuation.

Emphasis is usually marked with \emph{italic} text, which in an upright paragraph
reverses to roman when nested inside an italic context. Strong statements can be set in
\textbf{bold}, while \textit{explicit italic} and \textsc{small capitals} are useful for
names and terms that deserve a visual distinction without strong emphasis. Source code
and file names belong in \texttt{typewriter type}, and interface labels read best in
\textsf{sans serif}.

Font declarations act on the rest of the current group. The sentence {\bfseries this
phrase is set in bold by a declaration}, and {\itshape this one by an italic
declaration}, so that the surrounding text is unaffected. The older forms \bf and \it
work the same way inside braces, as in {\bf old-style bold} and {\it old-style italic}.

Accents are part of ordinary text. The word caf\'e takes an acute accent, the name
G\"odel takes an umlaut, and the phrase na\"\i ve reads as it should. Quotations are
typed as ``double quotes'' and the apostrophe in don't is a plain character, while
the dashes -- en and --- em -- have their own meanings.

Sentences that trail away end with \ldots{} and the typesetter spaces them as a unit.
The name of the system is written as \TeX{} in the running text and \LaTeX{} in the
document class that builds on it. Non-breaking spaces such as Figure~1 and Section~2
keep a label with its number, and a footnote gives a remark that would interrupt the
argument\footnote{A footnote is set at the foot of the page where its reference
appears, in a smaller size, and it is numbered in sequence with the other footnotes of
the document.}.

Explicit line breaks are useful in addresses and short verse:\\
the first line ends at the double backslash,\\
and the second line starts afresh.

Horizontal and vertical space can be inserted by hand. A paragraph that begins with
\noindent is not indented, and an explicit horizontal gap of \hspace{1em} separates
the two words in this sentence. Vertical space is more common between blocks, so the
next paragraph follows after \vspace{6pt} of extra room.

\noindent The paragraph above was started with the command \texttt{noindent}, so it
has no first-line indentation. The paragraphs that follow are indented as usual, which
is the default behaviour of the article class.

Ordinary paragraphs are separated by a blank line. The line breaking algorithm looks at
the whole paragraph before it commits to any break, which is why a change in one word
can move words several lines further down. The author rarely needs to intervene, but
the mechanisms exist, and the rest of this document uses them where they matter.

\subsection{Font declarations in context}\label{sec:fonts-7}

A declaration such as \textsf{bfseries} changes the font until the end of the enclosing
group. The next paragraph shows how groups limit the effect of such declarations, and
how the text returns to its original shape when the group closes.

{\bfseries Inside this group the text is bold, and a word such as \emph{this} still
changes to italic within bold text.} Outside the group the text returns to roman, as
the reader can see in this sentence. The same holds for {\itshape the slanted shape}
and for {\sffamily the sans serif family}, which all end where their braces end.

\subsection{Spacing and hyphenation}\label{sec:spacing-7}

Word spacing in a justified line is not constant. The stretch that a line receives is
shared among its interword glue, and the amount of stretch depends on how much room the
line has after its words are placed. A line with few words and a large gap is
noticeably loose, and a line with many words and little room may be tight. The
typesetter tries to keep all lines of a paragraph as even as possible, and it accepts
a slightly loose line when the alternative is a worse paragraph as a whole.

Long words that do not fit at the end of a line are hyphenated at points chosen from a
table of patterns. The table is built from word lists and the rules that govern where
a word may be divided, and it lets the typesetter break words such as typesetting,
hyphenation, algorithm and reproducibility at reasonable places. A hyphen at the end of
a line is a cost, and two such hyphens in a row are a larger cost, so the algorithm
avoids them where it can.

Some words must never be divided, and the author can say so by writing the word inside
a group that forbids breaks, or by joining the parts of a name with a tie. The name
Ludwig~van Beethoven keeps its two parts together, and a unit such as 12~pt does the
same for the number and the unit.

Sentences that end with a capital letter, such as the initials in J.~R.\ Smith, have a
different space after them, and the author must mark an end of sentence when the
capital is not the end of a sentence. The rule is simple, but it is easy to forget, and
the document below uses a backslash space in one case to show how it is written.

\section{Lists}\label{sec:lists-7}

Lists organise material that has no natural sentence structure. The three list
environments of the article class differ in how they mark their items and how they
number them. The first is the unordered list, which uses a bullet for each item.

\begin{itemize}
  \item The first item of a bulleted list is a plain sentence.
  \item The second item is longer and contains an \emph{emphasised} word, so that the
        wrapped line is indented to the same position as the text of the item.
  \item A third item may itself hold a sublist:
    \begin{itemize}
      \item the nested list uses a dash as its marker,
      \item and its items are indented further than the parent items.
    \end{itemize}
\end{itemize}

Numbered lists are used when the order of the items matters, for example in a
procedure. The following enumeration is nested two levels deep, and its first item
carries a label so that the text can refer to it later. The reference to that item
appears in Section~\ref{sec:tables-7}, which is written after this list.

\begin{enumerate}
  \item\label{item:first-7} Obtain the source files of the document and the packages that
        it loads.
  \item Compile the document once on the native path, which reads the auxiliary
        information of the previous run, if there is one.
        \begin{enumerate}
          \item The first sub-item describes a check that the author can make by hand.
          \item The second sub-item describes what happens when the check fails.
        \end{enumerate}
  \item Compare the result against the reference output and report any page that
        differs.
\end{enumerate}

The first step of the procedure is item~\ref{item:first-7} in the list above, and the
reader who wants the reasons behind the procedure will find them in
Section~\ref{sec:displays-7}.

Definition lists give a term followed by its explanation. The description environment
sets each term in bold and indents the explanation beneath it.

\begin{description}
  \item[Kerning] An adjustment of the space between two particular letters, such as
        the pair A and V, so that the shapes look evenly spaced.
  \item[Ligature] A single glyph that replaces a sequence of letters, such as the
        fi ligature in the word first.
  \item[Leading] The distance from one baseline to the next, measured in points.
\end{description}

The three kinds of list above are the core of most structured writing. An author who
needs a list inside a list simply nests the environment, and each level takes its own
label and its own indentation. The nesting depth in the article class is limited only
by the page, although a list nested more than three levels deep is usually a sign that
the material wants a table or a section instead.

Lists also interact with the page. A list is a sequence of paragraphs with extra space
between items, and the page builder may break the page between two items, as it may
between two paragraphs. The space at the top and bottom of a list is controlled by the
document class, so that a list opens and closes with the same visual distance as a
displayed equation.

\section{Displayed text}\label{sec:displays-7}

Some text is set apart from the body of the document to show that it is a quotation, a
centred title, or a block of literal material. Each environment has its own margins and
its own vertical spacing, and the environments below show the common cases.

\begin{quote}
A quotation environment sets its text with extra margin on both sides and without
paragraph indentation. It is intended for a short passage that the author wishes to
present as a separate block of prose, without the extra emphasis of a block quotation.
\end{quote}

\begin{quotation}
The quotation environment differs from the quote environment in that it indents the
first line of each of its paragraphs. A long passage that runs over several paragraphs
therefore keeps its internal structure, as a reader expects from a transcript of a
speech or a letter. This quotation has two paragraphs, so the second one can be seen to
start with an indent.

The second paragraph of the quotation shows the indentation of a continuing paragraph.
\end{quotation}

\begin{center}
Centred text on a line of its own\\
keeps each of its lines in the middle of the text block.
\end{center}

\begin{flushleft}
Flush left text keeps the left margin straight and lets the right margin
ragged, which suits material such as addresses and minutes of a meeting.
\end{flushleft}

\begin{flushright}
Flush right text keeps the right margin straight, which is seldom useful for
prose but is useful for a signature line.
\end{flushright}

Verbatim text is typeset exactly as it is typed, with spaces, line breaks and special
characters such as the percent sign and the braces kept as they stand:

\begin{verbatim}
  for i in 1..n {
      total += a[i] * b[i];   % dot product, not a comment
  }
\end{verbatim}

The following block uses a centring declaration inside a group, which keeps the
alignment local to the braces, and then a ragged right declaration on the next block.

{\centering A centred block made with a declaration, which ends at the closing
brace and leaves the surrounding text alone.\par}

{\raggedright A block whose lines are not stretched to the full measure, so the
right margin is ragged. The declaration ends with the group, and the next paragraph
returns to the usual justified setting of the text.\par}

Paragraphs of ordinary text after these blocks are justified again, as the reader can
see here.

\section{Tables}\label{sec:tables-7}

Tables set data in rows and columns. The tabular environment takes a column
specification, and each column is either left, centred or right aligned. Vertical bars
draw rules between columns. Section~\ref{sec:displays-7} explains the centring that surrounds each table in
this section.

\subsection{Basic alignment and rules}\label{sec:tables-basic-7}

\begin{center}
\begin{tabular}{l c r}
Method & Passes & Time (ms) \\
Greedy fill      & 1 & 0.4 \\
Knuth--Plass     & 2 & 3.1 \\
Optimal fit      & 3 & 4.7 \\
\end{tabular}
\end{center}

\begin{center}
\begin{tabular}{|l|c|r|}
Language & Typing & Speed \\
C        & static  & fast   \\
Python   & dynamic & slow   \\
Rust     & static  & fast   \\
\end{tabular}
\end{center}

\subsection{Padding and repeated columns}\label{sec:tables-advanced-7}

The column specification may fix the padding around a column with an empty at-sign
argument, which removes the space at one end of the table. Repeated columns are written
with a star argument, which expands to a run of identical columns.

\begin{center}
\begin{tabular}{@{}*{3}{c}@{}}
Left & Middle & Right \\
1 & 2 & 3 \\
4 & 5 & 6 \\
\end{tabular}
\end{center}

Each column keeps its own alignment in every row, so the numbers in the table below
line up on their right edges whatever the width of their labels.

\begin{center}
\begin{tabular}{l c r}
Run & Pages & Seconds \\
Alpha & 12 & 0.8 \\
Beta  & 29 & 1.9 \\
Gamma & 61 & 4.2 \\
\end{tabular}
\end{center}

Extra vertical space between rows is written after a row break with an optional length,
as in \texttt{[2pt]}. The table below uses that spacing to separate groups of rows.

\begin{center}
\begin{tabular}{l r}
Operation & Cost \\[2pt]
Read      & 3 \\
Write     & 7 \\[2pt]
Reference & 11 \\
\end{tabular}
\end{center}

Tables with many rows are common in reports, and the author who sets them by hand often
spends more time on alignment than on content. The tabular environment reduces that
work to a column specification and a list of rows. Each row ends with a double
backslash, and each cell is separated from its neighbour by an ampersand. The
alignment of each column is taken from the specification, so a change of one letter
moves the whole column without any other edit.

The tables so far show the alignment letters, the padding that the at-sign notation
controls and the repetition of a column pattern. Horizontal rules, partial rules and
cells that span several columns are set by the same machinery.

\subsection{Rules and spanning cells}\label{sec:tables-rules-7}
Table rules come in three kinds: a full rule across every column, a partial rule
across some of them, and a doubled rule made of two full rules in a row. A heading can
span all the columns of the table, and a single cell can span two of them.

\begin{center}
\begin{tabular}{|l|c|r|}
\hline
\multicolumn{3}{|c|}{\bfseries Measured costs} \\
\hline\hline
Phase & Calls & Time \\
\cline{2-3}
Scan    & 120 & 3.5 \\
Expand  & 940 & 12.0 \\
\multicolumn{2}{|l|}{Total, both phases} & 15.5 \\
\hline
\end{tabular}
\end{center}

\section{Floats}\label{sec:floats-7}

Floating material such as tables and figures is moved by the page builder to a position
where it fits. The placement specifier suggests where the float may go: \texttt{h} for
here, \texttt{t} for the top of a page, \texttt{b} for the bottom and \texttt{p} for a
page of floats. The author can give several letters at once, as in \texttt{[htbp]}, and
the page builder tries them in turn.

Table~\ref{tab:results-7} appears before the text that discusses it, so the reader meets
the reference first and the float later. The float is placed here by its own
specifier.

\begin{table}[h]
\centering
\begin{tabular}{l c r}
Input size & Lines & Pages \\
Short      & 40    & 1     \\
Medium     & 400   & 12    \\
Long       & 4000  & 120   \\
\end{tabular}
\caption{Pages produced for three input sizes.}\label{tab:results-7}
\end{table}

The results in Table~\ref{tab:results-7} grow with the input, and the page counts follow
the line counts closely. The next float is a figure with a caption long enough to wrap
onto a second line, and it is placed at the top of the page by its specifier.

\begin{figure}[t]
\centering
\mbox{\textsf{Input}\quad$\rightarrow$\quad\texttt{Parser}\quad$\rightarrow$\quad\textsf{Layout}\quad$\rightarrow$\quad\textsf{PDF}}
\caption{The stages of the native pipeline, from the characters of the source file to the
bytes of the output file, shown as a chain of boxes in the order in which the data flows
through them.}\label{fig:pipeline-7}
\end{figure}

Figure~\ref{fig:pipeline-7} shows the chain of stages that the native path runs in a
single pass. The figure caption is set as a paragraph because it runs over two lines.
A table placed at the bottom of a page is specified with \texttt{[b]}, as the next
example shows.

\begin{table}[b]
\centering
\begin{tabular}{c c}
Feature & State \\
Lists        & done \\
Tables       & done \\
\end{tabular}
\caption{Feature status at the time of writing.}
\end{table}

A float with both the top and the bottom permitted is written \texttt{[tb]}, and a float
that may go anywhere in the sequence is written \texttt{[htbp]}. The next figure uses
the latter specifier and a label that the text refers to from a later section.

\begin{figure}[htbp]
\centering
\mbox{\texttt{figure}} \quad \mbox{Text placed in a box}
\caption{A short caption.}\label{fig:box-7}
\end{figure}

Figure~\ref{fig:box-7} is referred to here, after the float itself. The references to the
tables and figures of this section are resolved in the second pass, which is the reason
that the reader sees the correct numbers even for references that come before their
floats.

A float is a block of material that the page builder may move away from the place where
it appears in the source. Its purpose is to keep large items, such as a table of
results or a diagram, from splitting across a page break or from leaving a large gap
at the foot of a page. The page builder reads the placement specifier of the float and
considers each candidate position in turn. If no candidate fits, the float waits for
the next page, and the material that follows it continues to fill the current page.

A caption belongs to its float and is numbered with it. The number of a table and the
number of a figure are kept in separate counters, so that a document may contain a
table 1 and a figure 1 without conflict. The label placed after the caption records the
number, and the reference in the text picks up the number from the label when the
document is compiled. This is the reason that a reference may come before its float in
the source, as the references to the tables above do.

\section{Mathematics}\label{sec:math-7}

Mathematical notation is set in math mode, which uses its own fonts and spacing rules.
Inline formulae are written between dollar signs. The quantity $\frac{a+b}{c-d}$ is a
fraction, and $\sqrt{x^2+1}$ is a square root. A superscript such as $x^{2}$ and a
subscript such as $a_{i}$ are common, and both can appear together, as in $x_{i}^{2}$.

Greek letters are typed by name, so $\alpha$, $\beta$ and $\gamma$ are available, and so
is the capital $\Gamma$. The sum $\sum_{k=1}^{n} k$ and the integral $\int_0^1 x\,dx$
appear inline with their limits set beside them, as is customary in running text. The
limit $\lim_{n\to\infty} \frac{1}{n} = 0$ also appears inline.

Sets are typeset with blackboard and calligraphic letters. The set of real numbers is
$\mathbb{R}$, and a function that bounds another is written $f(n) \in \mathcal{O}(n^2)$.
A vector is written in bold, as $\mathbf{v}$, and its length is $\|\mathbf{v}\|$. The
parentheses in $\left( \frac{x}{y} \right)$ grow with the fraction they enclose.

Displayed formulae are set on their own line, centred, with a larger style than inline
formulae. The formula below gives the sum of the first $n$ integers.
\[
  \sum_{k=1}^{n} k = \frac{n(n+1)}{2}
\]

Numbered equations carry a label, so that the text can refer to them by number. The
first equation of this section is the definition of the Gaussian integral, and the
second is the quadratic formula, which is used in Section~\ref{sec:lists-7}.

\begin{equation}\label{eq:gauss-7}
  \int_{-\infty}^{\infty} e^{-x^{2}}\,dx = \sqrt{\pi}
\end{equation}

\begin{equation}\label{eq:quadratic-7}
  x = \frac{-b \pm \sqrt{b^{2}-4ac}}{2a}
\end{equation}

Equation~\ref{eq:gauss-7} gives the Gaussian integral, and
Equation~\ref{eq:quadratic-7} gives the roots of a quadratic polynomial. The coefficients
$a$, $b$ and $c$ in Equation~\ref{eq:quadratic-7} are real numbers with $a \neq 0$.

A last display shows a sum over a product, with a subscript that is itself a formula:
\[
  \prod_{i=1}^{m} \left( 1 + \frac{1}{i^{2}} \right) \leq e^{\pi^{2}/6}
\]

Mathematical formulae are read as a sequence of symbols with a structure, and the layout
rules give each symbol a size and a position. Inline formulae are set in the text style,
which keeps the lines of a paragraph at a uniform height, while displayed formulae are
set in the display style, which allows larger operators and fractions with larger
numerators and denominators. The reader sees the difference in the sums and products
above, where the limits sit above and below the symbol in the display, and beside it in
the running text.

The quadratic formula in Equation~\ref{eq:quadratic-7} is a familiar example of a
fraction with a root in its numerator. The square root of the discriminant
$b^{2}-4ac$ is set with a radical sign that grows with its contents, and the
fraction line extends beneath the whole numerator. The parentheses of the next example
are sized in the same way, so that the reader can see the structure at a glance.

\subsection{Displayed derivations}\label{sec:derivations-7}

A derivation shows a chain of equalities, each following from the one before. Each line
of the chain is written as a separate display, and the text between the displays gives
the reason for each step. The first derivation below expands the square of a binomial,
and the second uses the result to evaluate a finite sum of squares.
\[
  (a+b)^{2} = a^{2} + 2ab + b^{2}
\]

Applying the same expansion to the terms of a sum gives a closed form for the sum of the
squares of the first $n$ integers, which is the second formula of this subsection:
\[
  \sum_{k=1}^{n} k^{2} = \frac{n(n+1)(2n+1)}{6}
\]

The formula holds for every positive integer $n$, and it can be proved by induction on
$n$, as the reader may check for the first few values. Each step of the induction
uses the identity above and the fact that the sum for $n+1$ is the sum for $n$ with one
more term added.

The third display is a limit that defines the constant $e$. It is set on its own line
because it is a result that the text refers to, and the reader may wish to compare it
with the bound given in the last display of the previous section:
\[
  e = \lim_{n\to\infty} \left( 1 + \frac{1}{n} \right)^{n}
\]

Derivations that are long are best broken into several displays, each short enough to
fit on one line of the text block, so that the reader can follow the argument without
turning the page.

\section{Bibliography and citations}\label{sec:refs-7}

The native path follows the algorithms of the typesetting system closely, and the
descriptions of those algorithms are found in the standard references. The line
breaking method used by the paragraph builder is described by Knuth and
Plass~\cite{knuth81-7}, while the page builder and the float placement are described in
the reference manual~\cite{knuth86-7}. The layout of mathematical formulae follows the
rules of the appendix on mathematics in that same manual. The chain of stages in
Figure~\ref{fig:pipeline-7} is a common design for compilers in general, as Aho and
colleagues describe~\cite{aho06-7}.

\begin{thebibliography}{9}
\bibitem{knuth81-7}
  Donald E.~Knuth and Michael F.~Plass.
  \emph{Breaking paragraphs into lines}.
  Software: Practice and Experience, 11(11), 1981.

\bibitem{knuth86-7}
  Donald E.~Knuth.
  \emph{The \TeX book}.
  Addison-Wesley, 1986.

\bibitem{aho06-7}
  Alfred V.~Aho, Monica S.~Lam, Ravi Sethi and Jeffrey D.~Ullman.
  \emph{Compilers: Principles, Techniques, and Tools}.
  Addison-Wesley, second edition, 2006.
\end{thebibliography}



\section{Running text}\label{sec:text-8}

Typesetting is the business of turning a stream of characters into a page. A document
author writes paragraphs of plain prose, and the typesetter decides where lines break,
how much space separates words, and where the page ends. The quality of the result
depends on many small decisions, most of which the author never sees. This section
collects the inline devices that an author reaches for while writing prose: emphasis,
font families, font shapes, accents and punctuation.

Emphasis is usually marked with \emph{italic} text, which in an upright paragraph
reverses to roman when nested inside an italic context. Strong statements can be set in
\textbf{bold}, while \textit{explicit italic} and \textsc{small capitals} are useful for
names and terms that deserve a visual distinction without strong emphasis. Source code
and file names belong in \texttt{typewriter type}, and interface labels read best in
\textsf{sans serif}.

Font declarations act on the rest of the current group. The sentence {\bfseries this
phrase is set in bold by a declaration}, and {\itshape this one by an italic
declaration}, so that the surrounding text is unaffected. The older forms \bf and \it
work the same way inside braces, as in {\bf old-style bold} and {\it old-style italic}.

Accents are part of ordinary text. The word caf\'e takes an acute accent, the name
G\"odel takes an umlaut, and the phrase na\"\i ve reads as it should. Quotations are
typed as ``double quotes'' and the apostrophe in don't is a plain character, while
the dashes -- en and --- em -- have their own meanings.

Sentences that trail away end with \ldots{} and the typesetter spaces them as a unit.
The name of the system is written as \TeX{} in the running text and \LaTeX{} in the
document class that builds on it. Non-breaking spaces such as Figure~1 and Section~2
keep a label with its number, and a footnote gives a remark that would interrupt the
argument\footnote{A footnote is set at the foot of the page where its reference
appears, in a smaller size, and it is numbered in sequence with the other footnotes of
the document.}.

Explicit line breaks are useful in addresses and short verse:\\
the first line ends at the double backslash,\\
and the second line starts afresh.

Horizontal and vertical space can be inserted by hand. A paragraph that begins with
\noindent is not indented, and an explicit horizontal gap of \hspace{1em} separates
the two words in this sentence. Vertical space is more common between blocks, so the
next paragraph follows after \vspace{6pt} of extra room.

\noindent The paragraph above was started with the command \texttt{noindent}, so it
has no first-line indentation. The paragraphs that follow are indented as usual, which
is the default behaviour of the article class.

Ordinary paragraphs are separated by a blank line. The line breaking algorithm looks at
the whole paragraph before it commits to any break, which is why a change in one word
can move words several lines further down. The author rarely needs to intervene, but
the mechanisms exist, and the rest of this document uses them where they matter.

\subsection{Font declarations in context}\label{sec:fonts-8}

A declaration such as \textsf{bfseries} changes the font until the end of the enclosing
group. The next paragraph shows how groups limit the effect of such declarations, and
how the text returns to its original shape when the group closes.

{\bfseries Inside this group the text is bold, and a word such as \emph{this} still
changes to italic within bold text.} Outside the group the text returns to roman, as
the reader can see in this sentence. The same holds for {\itshape the slanted shape}
and for {\sffamily the sans serif family}, which all end where their braces end.

\subsection{Spacing and hyphenation}\label{sec:spacing-8}

Word spacing in a justified line is not constant. The stretch that a line receives is
shared among its interword glue, and the amount of stretch depends on how much room the
line has after its words are placed. A line with few words and a large gap is
noticeably loose, and a line with many words and little room may be tight. The
typesetter tries to keep all lines of a paragraph as even as possible, and it accepts
a slightly loose line when the alternative is a worse paragraph as a whole.

Long words that do not fit at the end of a line are hyphenated at points chosen from a
table of patterns. The table is built from word lists and the rules that govern where
a word may be divided, and it lets the typesetter break words such as typesetting,
hyphenation, algorithm and reproducibility at reasonable places. A hyphen at the end of
a line is a cost, and two such hyphens in a row are a larger cost, so the algorithm
avoids them where it can.

Some words must never be divided, and the author can say so by writing the word inside
a group that forbids breaks, or by joining the parts of a name with a tie. The name
Ludwig~van Beethoven keeps its two parts together, and a unit such as 12~pt does the
same for the number and the unit.

Sentences that end with a capital letter, such as the initials in J.~R.\ Smith, have a
different space after them, and the author must mark an end of sentence when the
capital is not the end of a sentence. The rule is simple, but it is easy to forget, and
the document below uses a backslash space in one case to show how it is written.

\section{Lists}\label{sec:lists-8}

Lists organise material that has no natural sentence structure. The three list
environments of the article class differ in how they mark their items and how they
number them. The first is the unordered list, which uses a bullet for each item.

\begin{itemize}
  \item The first item of a bulleted list is a plain sentence.
  \item The second item is longer and contains an \emph{emphasised} word, so that the
        wrapped line is indented to the same position as the text of the item.
  \item A third item may itself hold a sublist:
    \begin{itemize}
      \item the nested list uses a dash as its marker,
      \item and its items are indented further than the parent items.
    \end{itemize}
\end{itemize}

Numbered lists are used when the order of the items matters, for example in a
procedure. The following enumeration is nested two levels deep, and its first item
carries a label so that the text can refer to it later. The reference to that item
appears in Section~\ref{sec:tables-8}, which is written after this list.

\begin{enumerate}
  \item\label{item:first-8} Obtain the source files of the document and the packages that
        it loads.
  \item Compile the document once on the native path, which reads the auxiliary
        information of the previous run, if there is one.
        \begin{enumerate}
          \item The first sub-item describes a check that the author can make by hand.
          \item The second sub-item describes what happens when the check fails.
        \end{enumerate}
  \item Compare the result against the reference output and report any page that
        differs.
\end{enumerate}

The first step of the procedure is item~\ref{item:first-8} in the list above, and the
reader who wants the reasons behind the procedure will find them in
Section~\ref{sec:displays-8}.

Definition lists give a term followed by its explanation. The description environment
sets each term in bold and indents the explanation beneath it.

\begin{description}
  \item[Kerning] An adjustment of the space between two particular letters, such as
        the pair A and V, so that the shapes look evenly spaced.
  \item[Ligature] A single glyph that replaces a sequence of letters, such as the
        fi ligature in the word first.
  \item[Leading] The distance from one baseline to the next, measured in points.
\end{description}

The three kinds of list above are the core of most structured writing. An author who
needs a list inside a list simply nests the environment, and each level takes its own
label and its own indentation. The nesting depth in the article class is limited only
by the page, although a list nested more than three levels deep is usually a sign that
the material wants a table or a section instead.

Lists also interact with the page. A list is a sequence of paragraphs with extra space
between items, and the page builder may break the page between two items, as it may
between two paragraphs. The space at the top and bottom of a list is controlled by the
document class, so that a list opens and closes with the same visual distance as a
displayed equation.

\section{Displayed text}\label{sec:displays-8}

Some text is set apart from the body of the document to show that it is a quotation, a
centred title, or a block of literal material. Each environment has its own margins and
its own vertical spacing, and the environments below show the common cases.

\begin{quote}
A quotation environment sets its text with extra margin on both sides and without
paragraph indentation. It is intended for a short passage that the author wishes to
present as a separate block of prose, without the extra emphasis of a block quotation.
\end{quote}

\begin{quotation}
The quotation environment differs from the quote environment in that it indents the
first line of each of its paragraphs. A long passage that runs over several paragraphs
therefore keeps its internal structure, as a reader expects from a transcript of a
speech or a letter. This quotation has two paragraphs, so the second one can be seen to
start with an indent.

The second paragraph of the quotation shows the indentation of a continuing paragraph.
\end{quotation}

\begin{center}
Centred text on a line of its own\\
keeps each of its lines in the middle of the text block.
\end{center}

\begin{flushleft}
Flush left text keeps the left margin straight and lets the right margin
ragged, which suits material such as addresses and minutes of a meeting.
\end{flushleft}

\begin{flushright}
Flush right text keeps the right margin straight, which is seldom useful for
prose but is useful for a signature line.
\end{flushright}

Verbatim text is typeset exactly as it is typed, with spaces, line breaks and special
characters such as the percent sign and the braces kept as they stand:

\begin{verbatim}
  for i in 1..n {
      total += a[i] * b[i];   % dot product, not a comment
  }
\end{verbatim}

The following block uses a centring declaration inside a group, which keeps the
alignment local to the braces, and then a ragged right declaration on the next block.

{\centering A centred block made with a declaration, which ends at the closing
brace and leaves the surrounding text alone.\par}

{\raggedright A block whose lines are not stretched to the full measure, so the
right margin is ragged. The declaration ends with the group, and the next paragraph
returns to the usual justified setting of the text.\par}

Paragraphs of ordinary text after these blocks are justified again, as the reader can
see here.

\section{Tables}\label{sec:tables-8}

Tables set data in rows and columns. The tabular environment takes a column
specification, and each column is either left, centred or right aligned. Vertical bars
draw rules between columns. Section~\ref{sec:displays-8} explains the centring that surrounds each table in
this section.

\subsection{Basic alignment and rules}\label{sec:tables-basic-8}

\begin{center}
\begin{tabular}{l c r}
Method & Passes & Time (ms) \\
Greedy fill      & 1 & 0.4 \\
Knuth--Plass     & 2 & 3.1 \\
Optimal fit      & 3 & 4.7 \\
\end{tabular}
\end{center}

\begin{center}
\begin{tabular}{|l|c|r|}
Language & Typing & Speed \\
C        & static  & fast   \\
Python   & dynamic & slow   \\
Rust     & static  & fast   \\
\end{tabular}
\end{center}

\subsection{Padding and repeated columns}\label{sec:tables-advanced-8}

The column specification may fix the padding around a column with an empty at-sign
argument, which removes the space at one end of the table. Repeated columns are written
with a star argument, which expands to a run of identical columns.

\begin{center}
\begin{tabular}{@{}*{3}{c}@{}}
Left & Middle & Right \\
1 & 2 & 3 \\
4 & 5 & 6 \\
\end{tabular}
\end{center}

Each column keeps its own alignment in every row, so the numbers in the table below
line up on their right edges whatever the width of their labels.

\begin{center}
\begin{tabular}{l c r}
Run & Pages & Seconds \\
Alpha & 12 & 0.8 \\
Beta  & 29 & 1.9 \\
Gamma & 61 & 4.2 \\
\end{tabular}
\end{center}

Extra vertical space between rows is written after a row break with an optional length,
as in \texttt{[2pt]}. The table below uses that spacing to separate groups of rows.

\begin{center}
\begin{tabular}{l r}
Operation & Cost \\[2pt]
Read      & 3 \\
Write     & 7 \\[2pt]
Reference & 11 \\
\end{tabular}
\end{center}

Tables with many rows are common in reports, and the author who sets them by hand often
spends more time on alignment than on content. The tabular environment reduces that
work to a column specification and a list of rows. Each row ends with a double
backslash, and each cell is separated from its neighbour by an ampersand. The
alignment of each column is taken from the specification, so a change of one letter
moves the whole column without any other edit.

The tables so far show the alignment letters, the padding that the at-sign notation
controls and the repetition of a column pattern. Horizontal rules, partial rules and
cells that span several columns are set by the same machinery.

\subsection{Rules and spanning cells}\label{sec:tables-rules-8}
Table rules come in three kinds: a full rule across every column, a partial rule
across some of them, and a doubled rule made of two full rules in a row. A heading can
span all the columns of the table, and a single cell can span two of them.

\begin{center}
\begin{tabular}{|l|c|r|}
\hline
\multicolumn{3}{|c|}{\bfseries Measured costs} \\
\hline\hline
Phase & Calls & Time \\
\cline{2-3}
Scan    & 120 & 3.5 \\
Expand  & 940 & 12.0 \\
\multicolumn{2}{|l|}{Total, both phases} & 15.5 \\
\hline
\end{tabular}
\end{center}

\section{Floats}\label{sec:floats-8}

Floating material such as tables and figures is moved by the page builder to a position
where it fits. The placement specifier suggests where the float may go: \texttt{h} for
here, \texttt{t} for the top of a page, \texttt{b} for the bottom and \texttt{p} for a
page of floats. The author can give several letters at once, as in \texttt{[htbp]}, and
the page builder tries them in turn.

Table~\ref{tab:results-8} appears before the text that discusses it, so the reader meets
the reference first and the float later. The float is placed here by its own
specifier.

\begin{table}[h]
\centering
\begin{tabular}{l c r}
Input size & Lines & Pages \\
Short      & 40    & 1     \\
Medium     & 400   & 12    \\
Long       & 4000  & 120   \\
\end{tabular}
\caption{Pages produced for three input sizes.}\label{tab:results-8}
\end{table}

The results in Table~\ref{tab:results-8} grow with the input, and the page counts follow
the line counts closely. The next float is a figure with a caption long enough to wrap
onto a second line, and it is placed at the top of the page by its specifier.

\begin{figure}[t]
\centering
\mbox{\textsf{Input}\quad$\rightarrow$\quad\texttt{Parser}\quad$\rightarrow$\quad\textsf{Layout}\quad$\rightarrow$\quad\textsf{PDF}}
\caption{The stages of the native pipeline, from the characters of the source file to the
bytes of the output file, shown as a chain of boxes in the order in which the data flows
through them.}\label{fig:pipeline-8}
\end{figure}

Figure~\ref{fig:pipeline-8} shows the chain of stages that the native path runs in a
single pass. The figure caption is set as a paragraph because it runs over two lines.
A table placed at the bottom of a page is specified with \texttt{[b]}, as the next
example shows.

\begin{table}[b]
\centering
\begin{tabular}{c c}
Feature & State \\
Lists        & done \\
Tables       & done \\
\end{tabular}
\caption{Feature status at the time of writing.}
\end{table}

A float with both the top and the bottom permitted is written \texttt{[tb]}, and a float
that may go anywhere in the sequence is written \texttt{[htbp]}. The next figure uses
the latter specifier and a label that the text refers to from a later section.

\begin{figure}[htbp]
\centering
\mbox{\texttt{figure}} \quad \mbox{Text placed in a box}
\caption{A short caption.}\label{fig:box-8}
\end{figure}

Figure~\ref{fig:box-8} is referred to here, after the float itself. The references to the
tables and figures of this section are resolved in the second pass, which is the reason
that the reader sees the correct numbers even for references that come before their
floats.

A float is a block of material that the page builder may move away from the place where
it appears in the source. Its purpose is to keep large items, such as a table of
results or a diagram, from splitting across a page break or from leaving a large gap
at the foot of a page. The page builder reads the placement specifier of the float and
considers each candidate position in turn. If no candidate fits, the float waits for
the next page, and the material that follows it continues to fill the current page.

A caption belongs to its float and is numbered with it. The number of a table and the
number of a figure are kept in separate counters, so that a document may contain a
table 1 and a figure 1 without conflict. The label placed after the caption records the
number, and the reference in the text picks up the number from the label when the
document is compiled. This is the reason that a reference may come before its float in
the source, as the references to the tables above do.

\section{Mathematics}\label{sec:math-8}

Mathematical notation is set in math mode, which uses its own fonts and spacing rules.
Inline formulae are written between dollar signs. The quantity $\frac{a+b}{c-d}$ is a
fraction, and $\sqrt{x^2+1}$ is a square root. A superscript such as $x^{2}$ and a
subscript such as $a_{i}$ are common, and both can appear together, as in $x_{i}^{2}$.

Greek letters are typed by name, so $\alpha$, $\beta$ and $\gamma$ are available, and so
is the capital $\Gamma$. The sum $\sum_{k=1}^{n} k$ and the integral $\int_0^1 x\,dx$
appear inline with their limits set beside them, as is customary in running text. The
limit $\lim_{n\to\infty} \frac{1}{n} = 0$ also appears inline.

Sets are typeset with blackboard and calligraphic letters. The set of real numbers is
$\mathbb{R}$, and a function that bounds another is written $f(n) \in \mathcal{O}(n^2)$.
A vector is written in bold, as $\mathbf{v}$, and its length is $\|\mathbf{v}\|$. The
parentheses in $\left( \frac{x}{y} \right)$ grow with the fraction they enclose.

Displayed formulae are set on their own line, centred, with a larger style than inline
formulae. The formula below gives the sum of the first $n$ integers.
\[
  \sum_{k=1}^{n} k = \frac{n(n+1)}{2}
\]

Numbered equations carry a label, so that the text can refer to them by number. The
first equation of this section is the definition of the Gaussian integral, and the
second is the quadratic formula, which is used in Section~\ref{sec:lists-8}.

\begin{equation}\label{eq:gauss-8}
  \int_{-\infty}^{\infty} e^{-x^{2}}\,dx = \sqrt{\pi}
\end{equation}

\begin{equation}\label{eq:quadratic-8}
  x = \frac{-b \pm \sqrt{b^{2}-4ac}}{2a}
\end{equation}

Equation~\ref{eq:gauss-8} gives the Gaussian integral, and
Equation~\ref{eq:quadratic-8} gives the roots of a quadratic polynomial. The coefficients
$a$, $b$ and $c$ in Equation~\ref{eq:quadratic-8} are real numbers with $a \neq 0$.

A last display shows a sum over a product, with a subscript that is itself a formula:
\[
  \prod_{i=1}^{m} \left( 1 + \frac{1}{i^{2}} \right) \leq e^{\pi^{2}/6}
\]

Mathematical formulae are read as a sequence of symbols with a structure, and the layout
rules give each symbol a size and a position. Inline formulae are set in the text style,
which keeps the lines of a paragraph at a uniform height, while displayed formulae are
set in the display style, which allows larger operators and fractions with larger
numerators and denominators. The reader sees the difference in the sums and products
above, where the limits sit above and below the symbol in the display, and beside it in
the running text.

The quadratic formula in Equation~\ref{eq:quadratic-8} is a familiar example of a
fraction with a root in its numerator. The square root of the discriminant
$b^{2}-4ac$ is set with a radical sign that grows with its contents, and the
fraction line extends beneath the whole numerator. The parentheses of the next example
are sized in the same way, so that the reader can see the structure at a glance.

\subsection{Displayed derivations}\label{sec:derivations-8}

A derivation shows a chain of equalities, each following from the one before. Each line
of the chain is written as a separate display, and the text between the displays gives
the reason for each step. The first derivation below expands the square of a binomial,
and the second uses the result to evaluate a finite sum of squares.
\[
  (a+b)^{2} = a^{2} + 2ab + b^{2}
\]

Applying the same expansion to the terms of a sum gives a closed form for the sum of the
squares of the first $n$ integers, which is the second formula of this subsection:
\[
  \sum_{k=1}^{n} k^{2} = \frac{n(n+1)(2n+1)}{6}
\]

The formula holds for every positive integer $n$, and it can be proved by induction on
$n$, as the reader may check for the first few values. Each step of the induction
uses the identity above and the fact that the sum for $n+1$ is the sum for $n$ with one
more term added.

The third display is a limit that defines the constant $e$. It is set on its own line
because it is a result that the text refers to, and the reader may wish to compare it
with the bound given in the last display of the previous section:
\[
  e = \lim_{n\to\infty} \left( 1 + \frac{1}{n} \right)^{n}
\]

Derivations that are long are best broken into several displays, each short enough to
fit on one line of the text block, so that the reader can follow the argument without
turning the page.

\section{Bibliography and citations}\label{sec:refs-8}

The native path follows the algorithms of the typesetting system closely, and the
descriptions of those algorithms are found in the standard references. The line
breaking method used by the paragraph builder is described by Knuth and
Plass~\cite{knuth81-8}, while the page builder and the float placement are described in
the reference manual~\cite{knuth86-8}. The layout of mathematical formulae follows the
rules of the appendix on mathematics in that same manual. The chain of stages in
Figure~\ref{fig:pipeline-8} is a common design for compilers in general, as Aho and
colleagues describe~\cite{aho06-8}.

\begin{thebibliography}{9}
\bibitem{knuth81-8}
  Donald E.~Knuth and Michael F.~Plass.
  \emph{Breaking paragraphs into lines}.
  Software: Practice and Experience, 11(11), 1981.

\bibitem{knuth86-8}
  Donald E.~Knuth.
  \emph{The \TeX book}.
  Addison-Wesley, 1986.

\bibitem{aho06-8}
  Alfred V.~Aho, Monica S.~Lam, Ravi Sethi and Jeffrey D.~Ullman.
  \emph{Compilers: Principles, Techniques, and Tools}.
  Addison-Wesley, second edition, 2006.
\end{thebibliography}



\section{Running text}\label{sec:text-9}

Typesetting is the business of turning a stream of characters into a page. A document
author writes paragraphs of plain prose, and the typesetter decides where lines break,
how much space separates words, and where the page ends. The quality of the result
depends on many small decisions, most of which the author never sees. This section
collects the inline devices that an author reaches for while writing prose: emphasis,
font families, font shapes, accents and punctuation.

Emphasis is usually marked with \emph{italic} text, which in an upright paragraph
reverses to roman when nested inside an italic context. Strong statements can be set in
\textbf{bold}, while \textit{explicit italic} and \textsc{small capitals} are useful for
names and terms that deserve a visual distinction without strong emphasis. Source code
and file names belong in \texttt{typewriter type}, and interface labels read best in
\textsf{sans serif}.

Font declarations act on the rest of the current group. The sentence {\bfseries this
phrase is set in bold by a declaration}, and {\itshape this one by an italic
declaration}, so that the surrounding text is unaffected. The older forms \bf and \it
work the same way inside braces, as in {\bf old-style bold} and {\it old-style italic}.

Accents are part of ordinary text. The word caf\'e takes an acute accent, the name
G\"odel takes an umlaut, and the phrase na\"\i ve reads as it should. Quotations are
typed as ``double quotes'' and the apostrophe in don't is a plain character, while
the dashes -- en and --- em -- have their own meanings.

Sentences that trail away end with \ldots{} and the typesetter spaces them as a unit.
The name of the system is written as \TeX{} in the running text and \LaTeX{} in the
document class that builds on it. Non-breaking spaces such as Figure~1 and Section~2
keep a label with its number, and a footnote gives a remark that would interrupt the
argument\footnote{A footnote is set at the foot of the page where its reference
appears, in a smaller size, and it is numbered in sequence with the other footnotes of
the document.}.

Explicit line breaks are useful in addresses and short verse:\\
the first line ends at the double backslash,\\
and the second line starts afresh.

Horizontal and vertical space can be inserted by hand. A paragraph that begins with
\noindent is not indented, and an explicit horizontal gap of \hspace{1em} separates
the two words in this sentence. Vertical space is more common between blocks, so the
next paragraph follows after \vspace{6pt} of extra room.

\noindent The paragraph above was started with the command \texttt{noindent}, so it
has no first-line indentation. The paragraphs that follow are indented as usual, which
is the default behaviour of the article class.

Ordinary paragraphs are separated by a blank line. The line breaking algorithm looks at
the whole paragraph before it commits to any break, which is why a change in one word
can move words several lines further down. The author rarely needs to intervene, but
the mechanisms exist, and the rest of this document uses them where they matter.

\subsection{Font declarations in context}\label{sec:fonts-9}

A declaration such as \textsf{bfseries} changes the font until the end of the enclosing
group. The next paragraph shows how groups limit the effect of such declarations, and
how the text returns to its original shape when the group closes.

{\bfseries Inside this group the text is bold, and a word such as \emph{this} still
changes to italic within bold text.} Outside the group the text returns to roman, as
the reader can see in this sentence. The same holds for {\itshape the slanted shape}
and for {\sffamily the sans serif family}, which all end where their braces end.

\subsection{Spacing and hyphenation}\label{sec:spacing-9}

Word spacing in a justified line is not constant. The stretch that a line receives is
shared among its interword glue, and the amount of stretch depends on how much room the
line has after its words are placed. A line with few words and a large gap is
noticeably loose, and a line with many words and little room may be tight. The
typesetter tries to keep all lines of a paragraph as even as possible, and it accepts
a slightly loose line when the alternative is a worse paragraph as a whole.

Long words that do not fit at the end of a line are hyphenated at points chosen from a
table of patterns. The table is built from word lists and the rules that govern where
a word may be divided, and it lets the typesetter break words such as typesetting,
hyphenation, algorithm and reproducibility at reasonable places. A hyphen at the end of
a line is a cost, and two such hyphens in a row are a larger cost, so the algorithm
avoids them where it can.

Some words must never be divided, and the author can say so by writing the word inside
a group that forbids breaks, or by joining the parts of a name with a tie. The name
Ludwig~van Beethoven keeps its two parts together, and a unit such as 12~pt does the
same for the number and the unit.

Sentences that end with a capital letter, such as the initials in J.~R.\ Smith, have a
different space after them, and the author must mark an end of sentence when the
capital is not the end of a sentence. The rule is simple, but it is easy to forget, and
the document below uses a backslash space in one case to show how it is written.

\section{Lists}\label{sec:lists-9}

Lists organise material that has no natural sentence structure. The three list
environments of the article class differ in how they mark their items and how they
number them. The first is the unordered list, which uses a bullet for each item.

\begin{itemize}
  \item The first item of a bulleted list is a plain sentence.
  \item The second item is longer and contains an \emph{emphasised} word, so that the
        wrapped line is indented to the same position as the text of the item.
  \item A third item may itself hold a sublist:
    \begin{itemize}
      \item the nested list uses a dash as its marker,
      \item and its items are indented further than the parent items.
    \end{itemize}
\end{itemize}

Numbered lists are used when the order of the items matters, for example in a
procedure. The following enumeration is nested two levels deep, and its first item
carries a label so that the text can refer to it later. The reference to that item
appears in Section~\ref{sec:tables-9}, which is written after this list.

\begin{enumerate}
  \item\label{item:first-9} Obtain the source files of the document and the packages that
        it loads.
  \item Compile the document once on the native path, which reads the auxiliary
        information of the previous run, if there is one.
        \begin{enumerate}
          \item The first sub-item describes a check that the author can make by hand.
          \item The second sub-item describes what happens when the check fails.
        \end{enumerate}
  \item Compare the result against the reference output and report any page that
        differs.
\end{enumerate}

The first step of the procedure is item~\ref{item:first-9} in the list above, and the
reader who wants the reasons behind the procedure will find them in
Section~\ref{sec:displays-9}.

Definition lists give a term followed by its explanation. The description environment
sets each term in bold and indents the explanation beneath it.

\begin{description}
  \item[Kerning] An adjustment of the space between two particular letters, such as
        the pair A and V, so that the shapes look evenly spaced.
  \item[Ligature] A single glyph that replaces a sequence of letters, such as the
        fi ligature in the word first.
  \item[Leading] The distance from one baseline to the next, measured in points.
\end{description}

The three kinds of list above are the core of most structured writing. An author who
needs a list inside a list simply nests the environment, and each level takes its own
label and its own indentation. The nesting depth in the article class is limited only
by the page, although a list nested more than three levels deep is usually a sign that
the material wants a table or a section instead.

Lists also interact with the page. A list is a sequence of paragraphs with extra space
between items, and the page builder may break the page between two items, as it may
between two paragraphs. The space at the top and bottom of a list is controlled by the
document class, so that a list opens and closes with the same visual distance as a
displayed equation.

\section{Displayed text}\label{sec:displays-9}

Some text is set apart from the body of the document to show that it is a quotation, a
centred title, or a block of literal material. Each environment has its own margins and
its own vertical spacing, and the environments below show the common cases.

\begin{quote}
A quotation environment sets its text with extra margin on both sides and without
paragraph indentation. It is intended for a short passage that the author wishes to
present as a separate block of prose, without the extra emphasis of a block quotation.
\end{quote}

\begin{quotation}
The quotation environment differs from the quote environment in that it indents the
first line of each of its paragraphs. A long passage that runs over several paragraphs
therefore keeps its internal structure, as a reader expects from a transcript of a
speech or a letter. This quotation has two paragraphs, so the second one can be seen to
start with an indent.

The second paragraph of the quotation shows the indentation of a continuing paragraph.
\end{quotation}

\begin{center}
Centred text on a line of its own\\
keeps each of its lines in the middle of the text block.
\end{center}

\begin{flushleft}
Flush left text keeps the left margin straight and lets the right margin
ragged, which suits material such as addresses and minutes of a meeting.
\end{flushleft}

\begin{flushright}
Flush right text keeps the right margin straight, which is seldom useful for
prose but is useful for a signature line.
\end{flushright}

Verbatim text is typeset exactly as it is typed, with spaces, line breaks and special
characters such as the percent sign and the braces kept as they stand:

\begin{verbatim}
  for i in 1..n {
      total += a[i] * b[i];   % dot product, not a comment
  }
\end{verbatim}

The following block uses a centring declaration inside a group, which keeps the
alignment local to the braces, and then a ragged right declaration on the next block.

{\centering A centred block made with a declaration, which ends at the closing
brace and leaves the surrounding text alone.\par}

{\raggedright A block whose lines are not stretched to the full measure, so the
right margin is ragged. The declaration ends with the group, and the next paragraph
returns to the usual justified setting of the text.\par}

Paragraphs of ordinary text after these blocks are justified again, as the reader can
see here.

\section{Tables}\label{sec:tables-9}

Tables set data in rows and columns. The tabular environment takes a column
specification, and each column is either left, centred or right aligned. Vertical bars
draw rules between columns. Section~\ref{sec:displays-9} explains the centring that surrounds each table in
this section.

\subsection{Basic alignment and rules}\label{sec:tables-basic-9}

\begin{center}
\begin{tabular}{l c r}
Method & Passes & Time (ms) \\
Greedy fill      & 1 & 0.4 \\
Knuth--Plass     & 2 & 3.1 \\
Optimal fit      & 3 & 4.7 \\
\end{tabular}
\end{center}

\begin{center}
\begin{tabular}{|l|c|r|}
Language & Typing & Speed \\
C        & static  & fast   \\
Python   & dynamic & slow   \\
Rust     & static  & fast   \\
\end{tabular}
\end{center}

\subsection{Padding and repeated columns}\label{sec:tables-advanced-9}

The column specification may fix the padding around a column with an empty at-sign
argument, which removes the space at one end of the table. Repeated columns are written
with a star argument, which expands to a run of identical columns.

\begin{center}
\begin{tabular}{@{}*{3}{c}@{}}
Left & Middle & Right \\
1 & 2 & 3 \\
4 & 5 & 6 \\
\end{tabular}
\end{center}

Each column keeps its own alignment in every row, so the numbers in the table below
line up on their right edges whatever the width of their labels.

\begin{center}
\begin{tabular}{l c r}
Run & Pages & Seconds \\
Alpha & 12 & 0.8 \\
Beta  & 29 & 1.9 \\
Gamma & 61 & 4.2 \\
\end{tabular}
\end{center}

Extra vertical space between rows is written after a row break with an optional length,
as in \texttt{[2pt]}. The table below uses that spacing to separate groups of rows.

\begin{center}
\begin{tabular}{l r}
Operation & Cost \\[2pt]
Read      & 3 \\
Write     & 7 \\[2pt]
Reference & 11 \\
\end{tabular}
\end{center}

Tables with many rows are common in reports, and the author who sets them by hand often
spends more time on alignment than on content. The tabular environment reduces that
work to a column specification and a list of rows. Each row ends with a double
backslash, and each cell is separated from its neighbour by an ampersand. The
alignment of each column is taken from the specification, so a change of one letter
moves the whole column without any other edit.

The tables so far show the alignment letters, the padding that the at-sign notation
controls and the repetition of a column pattern. Horizontal rules, partial rules and
cells that span several columns are set by the same machinery.

\subsection{Rules and spanning cells}\label{sec:tables-rules-9}
Table rules come in three kinds: a full rule across every column, a partial rule
across some of them, and a doubled rule made of two full rules in a row. A heading can
span all the columns of the table, and a single cell can span two of them.

\begin{center}
\begin{tabular}{|l|c|r|}
\hline
\multicolumn{3}{|c|}{\bfseries Measured costs} \\
\hline\hline
Phase & Calls & Time \\
\cline{2-3}
Scan    & 120 & 3.5 \\
Expand  & 940 & 12.0 \\
\multicolumn{2}{|l|}{Total, both phases} & 15.5 \\
\hline
\end{tabular}
\end{center}

\section{Floats}\label{sec:floats-9}

Floating material such as tables and figures is moved by the page builder to a position
where it fits. The placement specifier suggests where the float may go: \texttt{h} for
here, \texttt{t} for the top of a page, \texttt{b} for the bottom and \texttt{p} for a
page of floats. The author can give several letters at once, as in \texttt{[htbp]}, and
the page builder tries them in turn.

Table~\ref{tab:results-9} appears before the text that discusses it, so the reader meets
the reference first and the float later. The float is placed here by its own
specifier.

\begin{table}[h]
\centering
\begin{tabular}{l c r}
Input size & Lines & Pages \\
Short      & 40    & 1     \\
Medium     & 400   & 12    \\
Long       & 4000  & 120   \\
\end{tabular}
\caption{Pages produced for three input sizes.}\label{tab:results-9}
\end{table}

The results in Table~\ref{tab:results-9} grow with the input, and the page counts follow
the line counts closely. The next float is a figure with a caption long enough to wrap
onto a second line, and it is placed at the top of the page by its specifier.

\begin{figure}[t]
\centering
\mbox{\textsf{Input}\quad$\rightarrow$\quad\texttt{Parser}\quad$\rightarrow$\quad\textsf{Layout}\quad$\rightarrow$\quad\textsf{PDF}}
\caption{The stages of the native pipeline, from the characters of the source file to the
bytes of the output file, shown as a chain of boxes in the order in which the data flows
through them.}\label{fig:pipeline-9}
\end{figure}

Figure~\ref{fig:pipeline-9} shows the chain of stages that the native path runs in a
single pass. The figure caption is set as a paragraph because it runs over two lines.
A table placed at the bottom of a page is specified with \texttt{[b]}, as the next
example shows.

\begin{table}[b]
\centering
\begin{tabular}{c c}
Feature & State \\
Lists        & done \\
Tables       & done \\
\end{tabular}
\caption{Feature status at the time of writing.}
\end{table}

A float with both the top and the bottom permitted is written \texttt{[tb]}, and a float
that may go anywhere in the sequence is written \texttt{[htbp]}. The next figure uses
the latter specifier and a label that the text refers to from a later section.

\begin{figure}[htbp]
\centering
\mbox{\texttt{figure}} \quad \mbox{Text placed in a box}
\caption{A short caption.}\label{fig:box-9}
\end{figure}

Figure~\ref{fig:box-9} is referred to here, after the float itself. The references to the
tables and figures of this section are resolved in the second pass, which is the reason
that the reader sees the correct numbers even for references that come before their
floats.

A float is a block of material that the page builder may move away from the place where
it appears in the source. Its purpose is to keep large items, such as a table of
results or a diagram, from splitting across a page break or from leaving a large gap
at the foot of a page. The page builder reads the placement specifier of the float and
considers each candidate position in turn. If no candidate fits, the float waits for
the next page, and the material that follows it continues to fill the current page.

A caption belongs to its float and is numbered with it. The number of a table and the
number of a figure are kept in separate counters, so that a document may contain a
table 1 and a figure 1 without conflict. The label placed after the caption records the
number, and the reference in the text picks up the number from the label when the
document is compiled. This is the reason that a reference may come before its float in
the source, as the references to the tables above do.

\section{Mathematics}\label{sec:math-9}

Mathematical notation is set in math mode, which uses its own fonts and spacing rules.
Inline formulae are written between dollar signs. The quantity $\frac{a+b}{c-d}$ is a
fraction, and $\sqrt{x^2+1}$ is a square root. A superscript such as $x^{2}$ and a
subscript such as $a_{i}$ are common, and both can appear together, as in $x_{i}^{2}$.

Greek letters are typed by name, so $\alpha$, $\beta$ and $\gamma$ are available, and so
is the capital $\Gamma$. The sum $\sum_{k=1}^{n} k$ and the integral $\int_0^1 x\,dx$
appear inline with their limits set beside them, as is customary in running text. The
limit $\lim_{n\to\infty} \frac{1}{n} = 0$ also appears inline.

Sets are typeset with blackboard and calligraphic letters. The set of real numbers is
$\mathbb{R}$, and a function that bounds another is written $f(n) \in \mathcal{O}(n^2)$.
A vector is written in bold, as $\mathbf{v}$, and its length is $\|\mathbf{v}\|$. The
parentheses in $\left( \frac{x}{y} \right)$ grow with the fraction they enclose.

Displayed formulae are set on their own line, centred, with a larger style than inline
formulae. The formula below gives the sum of the first $n$ integers.
\[
  \sum_{k=1}^{n} k = \frac{n(n+1)}{2}
\]

Numbered equations carry a label, so that the text can refer to them by number. The
first equation of this section is the definition of the Gaussian integral, and the
second is the quadratic formula, which is used in Section~\ref{sec:lists-9}.

\begin{equation}\label{eq:gauss-9}
  \int_{-\infty}^{\infty} e^{-x^{2}}\,dx = \sqrt{\pi}
\end{equation}

\begin{equation}\label{eq:quadratic-9}
  x = \frac{-b \pm \sqrt{b^{2}-4ac}}{2a}
\end{equation}

Equation~\ref{eq:gauss-9} gives the Gaussian integral, and
Equation~\ref{eq:quadratic-9} gives the roots of a quadratic polynomial. The coefficients
$a$, $b$ and $c$ in Equation~\ref{eq:quadratic-9} are real numbers with $a \neq 0$.

A last display shows a sum over a product, with a subscript that is itself a formula:
\[
  \prod_{i=1}^{m} \left( 1 + \frac{1}{i^{2}} \right) \leq e^{\pi^{2}/6}
\]

Mathematical formulae are read as a sequence of symbols with a structure, and the layout
rules give each symbol a size and a position. Inline formulae are set in the text style,
which keeps the lines of a paragraph at a uniform height, while displayed formulae are
set in the display style, which allows larger operators and fractions with larger
numerators and denominators. The reader sees the difference in the sums and products
above, where the limits sit above and below the symbol in the display, and beside it in
the running text.

The quadratic formula in Equation~\ref{eq:quadratic-9} is a familiar example of a
fraction with a root in its numerator. The square root of the discriminant
$b^{2}-4ac$ is set with a radical sign that grows with its contents, and the
fraction line extends beneath the whole numerator. The parentheses of the next example
are sized in the same way, so that the reader can see the structure at a glance.

\subsection{Displayed derivations}\label{sec:derivations-9}

A derivation shows a chain of equalities, each following from the one before. Each line
of the chain is written as a separate display, and the text between the displays gives
the reason for each step. The first derivation below expands the square of a binomial,
and the second uses the result to evaluate a finite sum of squares.
\[
  (a+b)^{2} = a^{2} + 2ab + b^{2}
\]

Applying the same expansion to the terms of a sum gives a closed form for the sum of the
squares of the first $n$ integers, which is the second formula of this subsection:
\[
  \sum_{k=1}^{n} k^{2} = \frac{n(n+1)(2n+1)}{6}
\]

The formula holds for every positive integer $n$, and it can be proved by induction on
$n$, as the reader may check for the first few values. Each step of the induction
uses the identity above and the fact that the sum for $n+1$ is the sum for $n$ with one
more term added.

The third display is a limit that defines the constant $e$. It is set on its own line
because it is a result that the text refers to, and the reader may wish to compare it
with the bound given in the last display of the previous section:
\[
  e = \lim_{n\to\infty} \left( 1 + \frac{1}{n} \right)^{n}
\]

Derivations that are long are best broken into several displays, each short enough to
fit on one line of the text block, so that the reader can follow the argument without
turning the page.

\section{Bibliography and citations}\label{sec:refs-9}

The native path follows the algorithms of the typesetting system closely, and the
descriptions of those algorithms are found in the standard references. The line
breaking method used by the paragraph builder is described by Knuth and
Plass~\cite{knuth81-9}, while the page builder and the float placement are described in
the reference manual~\cite{knuth86-9}. The layout of mathematical formulae follows the
rules of the appendix on mathematics in that same manual. The chain of stages in
Figure~\ref{fig:pipeline-9} is a common design for compilers in general, as Aho and
colleagues describe~\cite{aho06-9}.

\begin{thebibliography}{9}
\bibitem{knuth81-9}
  Donald E.~Knuth and Michael F.~Plass.
  \emph{Breaking paragraphs into lines}.
  Software: Practice and Experience, 11(11), 1981.

\bibitem{knuth86-9}
  Donald E.~Knuth.
  \emph{The \TeX book}.
  Addison-Wesley, 1986.

\bibitem{aho06-9}
  Alfred V.~Aho, Monica S.~Lam, Ravi Sethi and Jeffrey D.~Ullman.
  \emph{Compilers: Principles, Techniques, and Tools}.
  Addison-Wesley, second edition, 2006.
\end{thebibliography}

\end{document}
